% Pertemuan 1 Praktikum Pemrograman (IF25-11002). Dihasilkan oleh tools/build.py dari tools/konten/p01.py.
% Kompilasi: xelatex P01.tex (dua kali). Catatan: xelatex -jobname=P01-catatan "\def\catatan{1}% Pertemuan 1 Praktikum Pemrograman (IF25-11002). Dihasilkan oleh tools/build.py dari tools/konten/p01.py.
% Kompilasi: xelatex P01.tex (dua kali). Catatan: xelatex -jobname=P01-catatan "\def\catatan{1}% Pertemuan 1 Praktikum Pemrograman (IF25-11002). Dihasilkan oleh tools/build.py dari tools/konten/p01.py.
% Kompilasi: xelatex P01.tex (dua kali). Catatan: xelatex -jobname=P01-catatan "\def\catatan{1}% Pertemuan 1 Praktikum Pemrograman (IF25-11002). Dihasilkan oleh tools/build.py dari tools/konten/p01.py.
% Kompilasi: xelatex P01.tex (dua kali). Catatan: xelatex -jobname=P01-catatan "\def\catatan{1}\input{P01.tex}"
\documentclass[t]{beamer}
\usetheme{ITERA}
\input{catatan-preamble.tex}
\renewcommand\ITERAmeeting{Pertemuan 1}
\begin{document}
\SlideJudul{IF25-11002 PRAKTIKUM PEMROGRAMAN}{Pertemuan 1\\Hello World dan Struktur Program C++}{Program pertama, cara compiler membacanya, dan cara membaca pesan error-nya}{Tim Pengajar}{Tahun 2026. Sejajar dengan Pertemuan 1 Algoritma Pemrograman: Computational Thinking.}
\note{Sapa kelas dan perkenalkan tim pengajar. Hari ini semua orang pulang dengan program yang benar-benar jalan, termasuk yang belum pernah menulis kode.}

\begin{frame}{Dua mata kuliah, satu minggu}
\framesubtitle{Sebelum mulai}
Topik minggu ini sama di dua mata kuliah, dengan peran yang berbeda.
\vspace{10pt}

\begin{columns}[T]
\begin{column}{0.47\textwidth}
\begin{Blok}{Algoritma: Computational Thinking}
Memecah masalah, menulis pseudocode dan flowchart, menelusuri dengan trace table. Minggu ini: computational thinking.
\end{Blok}
\end{column}
\begin{column}{0.47\textwidth}
\begin{BlokGelap}{Praktikum: Hello World}
Menuliskan langkah itu dalam C++, menjalankannya, dan memperbaikinya saat salah. Minggu ini: program pertama.
\end{BlokGelap}
\end{column}
\end{columns}
\note{Tanyakan apa yang sudah dibahas di Algoritma minggu ini. Gunakan jawabannya sebagai jembatan ke praktikum.}
\end{frame}

\begin{frame}{Saat pulang nanti, kalian bisa}
\framesubtitle{Target hari ini}
\begin{columns}[T]
\begin{column}{0.47\textwidth}
\begin{Langkah}{1}{Menerapkan}
Menulis, mengompilasi, dan menjalankan program C++ yang menampilkan keluaran terformat dengan \texttt{cout}, \texttt{endl}, dan escape sequence

{\footnotesize\color{ITERAInkMuted}Sub-CPMK 1.1, CPMK0607}
\end{Langkah}
\end{column}
\begin{column}{0.47\textwidth}
\begin{Langkah}{2}{Menjelaskan}
Menjelaskan peran setiap bagian program C++ dan alur dari kode sumber sampai program berjalan, termasuk membaca pesan error kompilasi

{\footnotesize\color{ITERAInkMuted}Sub-CPMK 1.2, CPMK0608}
\end{Langkah}
\end{column}
\end{columns}
\vspace{8pt}
\begin{Catatan}
Dua kemampuan ini dinilai sama berat di Exit Ticket, tugas, UTS, dan UAS.
\end{Catatan}
\note{Bacakan target dengan pelan. Kata kuncinya menerapkan dan menjelaskan.}
\end{frame}

\begin{frame}{Cara kalian dinilai}
\framesubtitle{Penilaian semester}
\small
\begin{tabular}{@{}>{\raggedright\arraybackslash}p{251pt}>{\raggedright\arraybackslash}p{116pt}>{\raggedright\arraybackslash}p{242pt}>{\raggedright\arraybackslash}p{242pt}@{}}
\kepala{Komponen} & \kepala{Bobot} & \kepala{Bagian A, CPMK0607} & \kepala{Bagian B, CPMK0608}\\
Exit Ticket, 14 kali & 25\% & 12,5 & 12,5\\
\barisgenap Tugas, 12 kali & 25\% & 12,5 & 12,5\\
UTS, pertemuan 8 & 25\% & 12,5 & 12,5\\
\barisgenap UAS, pertemuan 16 & 25\% & 12,5 & 12,5\\
\textbf{Total} & \textbf{100\%} & \textbf{50} & \textbf{50}\\
\garisdasar
\end{tabular}
\vspace{10pt}

Setiap instrumen dibagi rata: Bagian A soal kode, Bagian B soal tracing dan penjelasan. Nilai minimum lulus tiap CPMK adalah 50.\note{Karena setiap instrumen selalu 50:50, kontribusi tiap CPMK pasti 50 poin apa pun bobot komponennya. Sebutkan rubrik empat level yang dipakai di semua tugas.}
\end{frame}

\begin{frame}{Alur 150 menit setiap pertemuan}
\framesubtitle{Siklus yang sama setiap minggu}
\vspace{20pt}
\begin{center}
\begin{tikzpicture}[node distance=14pt]
\node[fase=ITERAGold50,text=ITERAInk] (f0) {{\ITERADisplay\bfseries\fontsize{26pt}{30pt}\selectfont 15'}\\[4pt]{\bfseries Review}\\[2pt]{\footnotesize masalah pembuka}};
\node[fase=ITERAGold50,text=ITERAInk,right=of f0] (f1) {{\ITERADisplay\bfseries\fontsize{26pt}{30pt}\selectfont 25'}\\[4pt]{\bfseries Struggle}\\[2pt]{\footnotesize coba sendiri dulu}};
\draw[-{Stealth[length=7pt]},line width=1.5pt,ITERAStructure] (f0) -- (f1);
\node[fase=ITERAGold,text=ITERAInk,right=of f1] (f2) {{\ITERADisplay\bfseries\fontsize{26pt}{30pt}\selectfont 40'}\\[4pt]{\bfseries Materi}\\[2pt]{\footnotesize tebak lalu buktikan}};
\draw[-{Stealth[length=7pt]},line width=1.5pt,ITERAStructure] (f1) -- (f2);
\node[fase=ITERAGold50,text=ITERAInk,right=of f2] (f3) {{\ITERADisplay\bfseries\fontsize{26pt}{30pt}\selectfont 45'}\\[4pt]{\bfseries Hands-on}\\[2pt]{\footnotesize latihan bertingkat}};
\draw[-{Stealth[length=7pt]},line width=1.5pt,ITERAStructure] (f2) -- (f3);
\node[fase=ITERAGround,text=ITERAOnGround,right=of f3] (f4) {{\ITERADisplay\bfseries\fontsize{26pt}{30pt}\selectfont 25'}\\[4pt]{\bfseries Exit Ticket}\\[2pt]{\footnotesize sendiri, tanpa bantuan}};
\draw[-{Stealth[length=7pt]},line width=1.5pt,ITERAStructure] (f3) -- (f4);
\end{tikzpicture}
\end{center}
\vspace{16pt}
Urutannya sengaja: kalian mencoba dulu sebelum dijelaskan, dan menebak dulu sebelum melihat keluaran.\note{Ini filosofi problem-first dari RPS. Exit Ticket dimulai paling lambat menit ke-125.}
\end{frame}

\Bagian{1}{Masalah pembuka}{Kartu nama di layar}
\note{Masalah pembuka 10 menit.}

\begin{frame}[fragile]{Kartu nama di layar}
\framesubtitle{Masalah minggu ini}
\begin{columns}[T]
\begin{column}{0.55\textwidth}
{\small Bayangkan panitia PKKMB meminta setiap mahasiswa baru menampilkan kartu nama di layar komputer lab. Isinya nama, NIM, dan program studi, dengan garis pembatas di atas dan bawah.

\vspace{6pt}
Kalian tidak boleh mengetik teks itu di Word atau Notepad. Teks harus muncul karena sebuah program menyuruh komputer menampilkannya.\par}
\vspace{6pt}
\begin{Catatan}
\textbf{Diskusi 2 menit.} Kalau komputer hanya mengerti instruksi yang sangat persis, instruksi seperti apa yang perlu kalian tulis? Bagaimana komputer tahu kapan harus pindah baris?
\end{Catatan}
\end{column}
\begin{column}{0.41\textwidth}
\begin{Keluaran}[]
========================
    KARTU MAHASISWA
========================
Nama  : Rani Puspita
NIM   : 126140001
Prodi : Teknik Informatika
========================
\end{Keluaran}
\end{column}
\end{columns}
\note{Tampung beberapa jawaban tanpa menilai benar salah.}
\end{frame}

\begin{frame}{Dari langkah ke instruksi}
\framesubtitle{Sambungan dengan Algoritma}
\begin{columns}[T]
\begin{column}{0.31\textwidth}
\begin{Langkah}{1}{Pecah}
Kartu terdiri dari tujuh baris teks.
\end{Langkah}
\end{column}
\begin{column}{0.31\textwidth}
\begin{Langkah}{2}{Urutkan}
Tampilkan baris pertama, pindah baris, lalu baris kedua, dan seterusnya.
\end{Langkah}
\end{column}
\begin{column}{0.31\textwidth}
\begin{Langkah}{3}{Terjemahkan}
Tulis setiap langkah dalam bahasa yang dimengerti komputer.
\end{Langkah}
\end{column}
\end{columns}
\vspace{10pt}
Langkah 1 dan 2 dilatih di Algoritma. Langkah 3 kita kerjakan hari ini dengan C++.\note{Hubungkan eksplisit ke Algoritma minggu ini.}
\end{frame}

\Bagian{2}{Coba dulu, baru dijelaskan}{25 menit eksperimen di GDB Online}
\note{Struggle 25 menit. Berkeliling, jangan menjelaskan materi dulu.}

\begin{frame}{Misi kalian}
\framesubtitle{Struggle, 25 menit}
\begin{columns}[T]
\begin{column}{0.31\textwidth}
\begin{Langkah}{1}{Buka}
onlinegdb.com, pilih bahasa C++ di kanan atas.
\end{Langkah}
\end{column}
\begin{column}{0.31\textwidth}
\begin{Langkah}{2}{Ubah}
Ganti teks Hello World menjadi Halo, nama saya ..., lalu klik Run.
\end{Langkah}
\end{column}
\begin{column}{0.31\textwidth}
\begin{Langkah}{3}{Kembangkan}
Tampilkan dua baris: nama di baris pertama, NIM di baris kedua.
\end{Langkah}
\end{column}
\end{columns}
\vspace{8pt}
\begin{columns}[T]
\begin{column}{0.47\textwidth}
\begin{Blok}{Boleh}
Bertanya ke teman, mengubah apa saja, sengaja membuat error.
\end{Blok}
\end{column}
\begin{column}{0.47\textwidth}
\begin{BlokAlert}{Belum dulu}
Mencari jawaban jadi di internet atau AI.
\end{BlokAlert}
\end{column}
\end{columns}
\note{Catat kesalahan yang sering muncul untuk dibahas di debrief.}
\end{frame}

\begin{frame}{Apa yang kalian temukan?}
\framesubtitle{Debrief struggle}
\begin{columns}[T]
\begin{column}{0.31\textwidth}
\begin{BlokGelap}{Pertanyaan 1}
Apa yang terjadi kalau satu titik koma dihapus?
\end{BlokGelap}
\end{column}
\begin{column}{0.31\textwidth}
\begin{BlokGelap}{Pertanyaan 2}
Bagaimana cara kalian membuat dua baris?
\end{BlokGelap}
\end{column}
\begin{column}{0.31\textwidth}
\begin{BlokGelap}{Pertanyaan 3}
Baris mana yang menurut kalian tidak boleh dihapus?
\end{BlokGelap}
\end{column}
\end{columns}
\vspace{10pt}
Jawaban kalian akan kita uji di bagian berikutnya.\note{Tulis jawaban di papan; buktikan atau patahkan di bagian materi.}
\end{frame}

\Bagian{3}{Membedah program pertama}{Dari kode sampai keluaran, baris demi baris}
\note{Materi 40 menit. Tunjukkan setiap contoh langsung di GDB Online.}

\begin{frame}{Apa itu program, dan mengapa harus dikompilasi}
\framesubtitle{Kompilasi}
\small
\begin{tabular}{@{}>{\raggedright\arraybackslash}p{194pt}>{\raggedright\arraybackslash}p{681pt}@{}}
\kepala{Tahap} & \kepala{Yang terjadi}\\
1. Kode sumber & Kalian menulis instruksi C++ di file berakhiran \texttt{.cpp}.\\
\barisgenap 2. Kompilasi & Compiler (misalnya \texttt{g++}) memeriksa tata bahasa lalu menerjemahkan ke bahasa mesin.\\
3. Executable & Hasil terjemahan berupa program yang siap dijalankan.\\
\barisgenap 4. Jalankan & Program menampilkan keluaran di layar.\\
\garisdasar
\end{tabular}
\end{frame}

\begin{frame}[fragile]{Membedah Hello World baris demi baris}
\framesubtitle{Hello World}
\begin{columns}[T]
\begin{column}{0.55\textwidth}
\begin{Kode}[]{hello.cpp}
#include <iostream>
using namespace std;

int main() {
    cout << "Hello World!" << endl;
    return 0;
}
\end{Kode}
\end{column}
\begin{column}{0.41\textwidth}
\only<1>{\begin{TebakDulu}
{\ITERAKickerFont\fontsize{10pt}{12pt}\selectfont\color{ITERAAlert}TEBAK DULU}\par\vspace{4pt}
Sebelum dijalankan, tulis keluarannya di kertas.
\end{TebakDulu}
}\only<2>{\KeluaranFile[]{keluaran/P01/k001.txt}
}
\end{column}
\end{columns}
\note<1>{Minta mahasiswa menebak keluaran sebelum membuka slide berikutnya. Tanyakan satu atau dua tebakan yang berbeda, lalu buktikan.}
\end{frame}

\begin{frame}[fragile]{Menyusun keluaran dengan cout, endl, dan \textbackslash{}n}
\framesubtitle{cout}
\begin{columns}[T]
\begin{column}{0.55\textwidth}
\begin{Kode}[]{rangkai.cpp}
cout << "Hai" << "Dunia";
cout << "Hai " << "Dunia" << endl;
cout << "Baris 1" << endl << "Baris 2\n";
cout << "Baris 3\nBaris 4" << endl;
\end{Kode}
\end{column}
\begin{column}{0.41\textwidth}
\only<1>{\begin{TebakDulu}
{\ITERAKickerFont\fontsize{10pt}{12pt}\selectfont\color{ITERAAlert}TEBAK DULU}\par\vspace{4pt}
Sebelum dijalankan, tulis keluarannya di kertas.
\end{TebakDulu}
}\only<2>{\KeluaranFile[]{keluaran/P01/k002.txt}
}
\end{column}
\end{columns}
\note<1>{Minta mahasiswa menebak keluaran sebelum membuka slide berikutnya. Tanyakan satu atau dua tebakan yang berbeda, lalu buktikan.}
\end{frame}

\begin{frame}{Escape sequence untuk karakter khusus}
\framesubtitle{Escape sequence}
\small
\begin{tabular}{@{}>{\raggedright\arraybackslash}p{113pt}>{\raggedright\arraybackslash}p{246pt}>{\raggedright\arraybackslash}p{246pt}>{\raggedright\arraybackslash}p{246pt}@{}}
\kepala{Escape} & \kepala{Arti} & \kepala{Ditulis} & \kepala{Tampil}\\
\texttt{\textbackslash{}n} & baris baru & \texttt{"A\textbackslash{}nB"} & A dan B di dua baris\\
\barisgenap \texttt{\textbackslash{}t} & tab & \texttt{"A\textbackslash{}tB"} & A, jarak tab, B\\
\texttt{\textbackslash{}\textbackslash{}} & garis miring terbalik & \texttt{"C:\textbackslash{}\textbackslash{}data"} & \texttt{C:\textbackslash{}data}\\
\barisgenap \texttt{\textbackslash{}"} & kutip ganda & \texttt{"Dia \textbackslash{}"hai\textbackslash{}""} & \texttt{Dia "hai"}\\
\garisdasar
\end{tabular}
\end{frame}

\begin{frame}[fragile]{Tebak keluarannya}
\framesubtitle{Latihan menjelaskan, 3 menit}
\begin{columns}[T]
\begin{column}{0.56\textwidth}
\begin{Kode}[]{tebak.cpp}
cout << "A" << "B" << endl;
cout << "C\tD\n" << "E";
cout << "\"F\"" << endl;
\end{Kode}
\end{column}
\begin{column}{0.40\textwidth}
\begin{Catatan}
Tulis keluarannya di kertas, karakter demi karakter. Jangan dijalankan dulu.
\end{Catatan}
\end{column}
\end{columns}
\note{Contoh soal Bagian B. Beri waktu 3 menit sebelum slide jawaban.}
\end{frame}

\begin{frame}[fragile]{Tebak keluarannya: jawaban}
\framesubtitle{Latihan menjelaskan}
\begin{columns}[T]
\begin{column}{0.56\textwidth}
\begin{Kode}[]{tebak.cpp}
cout << "A" << "B" << endl;
cout << "C\tD\n" << "E";
cout << "\"F\"" << endl;
\end{Kode}
\end{column}
\begin{column}{0.40\textwidth}
\begin{Keluaran}[]
AB
C       D
E"F"
\end{Keluaran}
\begin{Catatan}
Baris ketiga: \texttt{E} dan \texttt{"F"} menyatu karena tidak ada pindah baris di antaranya. Jarak antara C dan D adalah satu tab.
\end{Catatan}
\end{column}
\end{columns}
\note{Tanyakan jawaban yang paling sering salah, lalu telusuri bersama.}
\end{frame}

\begin{frame}{Error yang sering muncul minggu ini}
\framesubtitle{Pesan diambil langsung dari g++}
\small
\begin{tabular}{@{}>{\raggedright\arraybackslash}p{208pt}>{\raggedright\arraybackslash}p{256pt}>{\raggedright\arraybackslash}p{398pt}@{}}
\kepala{Kesalahan} & \kepala{Contoh} & \kepala{Pesan pertama dari g++}\\
Lupa titik koma & \texttt{cout << "Hai"} & \texttt{error: expected ';' before 'return'}\\
\barisgenap Salah ketik nama & \texttt{count << "Hai";} & \texttt{error: 'count' was not declared in this scope}\\
Lupa \#include & \texttt{tanpa baris 1} & \texttt{error: 'cout' was not declared in this scope}\\
\barisgenap Kutip tidak ditutup & \texttt{cout << "Hai;} & \texttt{error: missing terminating " character}\\
Garis miring tunggal & \texttt{"C:\textbackslash{}data"} & \texttt{warning: unknown escape sequence: '\textbackslash{}d'}\\
\garisdasar
\end{tabular}
\vspace{10pt}

{\small Pesan di tabel ini dihasilkan dengan g++ dan opsi -Wall, sama seperti di cek.sh.}
\note{Minta mahasiswa memotret slide ini.}
\end{frame}

\Bagian{4}{Hands-on bertingkat}{45 menit, dari dasar sampai tantangan}
\note{Mahasiswa membuka folder 02 Hands-on/P1 - Hello World - Hands-on.}

\begin{frame}{Peta hands-on}
\framesubtitle{Folder 02 Hands-on/P1 - Hello World - Hands-on}
\small
\begin{tabular}{@{}>{\raggedright\arraybackslash}p{36pt}>{\raggedright\arraybackslash}p{267pt}>{\raggedright\arraybackslash}p{110pt}>{\raggedright\arraybackslash}p{83pt}>{\raggedright\arraybackslash}p{341pt}@{}}
\kepala{No} & \kepala{File} & \kepala{Tingkat} & \kepala{Waktu} & \kepala{Yang dilatih}\\
1 & \texttt{handson1-biodata.cpp} & Dasar & 10' & \texttt{cout}, \texttt{endl}, tiga baris keluaran\\
\barisgenap 2 & \texttt{handson2-kartu.cpp} & Menengah & 15' & garis pembatas, perataan, \texttt{\textbackslash{}t}, \texttt{\textbackslash{}"}\\
3 & \texttt{handson3-perbaiki.cpp} & Tantangan & 20' & lima kesalahan, perbaiki dan jelaskan\\
\barisgenap + & \texttt{bonus-rumah.cpp} & Pengayaan & sisa & \texttt{\textbackslash{}\textbackslash{}} di dalam ASCII art\\
\garisdasar
\end{tabular}
\vspace{10pt}

Kerjakan berurutan. Cocokkan keluaran dengan folder \texttt{expected/}; latihan yang membaca masukan memakai data di folder \texttt{input/}.\note{Saat berkeliling, minta mahasiswa yang mengerjakan latihan tantangan menjelaskan blok PENJELASAN mereka.}
\end{frame}

\begin{frame}{Cara mengecek dan aturannya}
\framesubtitle{Hands-on}
\begin{columns}[T]
\begin{column}{0.47\textwidth}
\begin{Blok}{Di GDB Online}
Salin isi file latihan, lengkapi TODO, klik Run. Bila program membaca masukan, ketik isi file di \texttt{input/}. Bandingkan dengan \texttt{expected/}.
\end{Blok}
\end{column}
\begin{column}{0.47\textwidth}
\begin{BlokGelap}{Di komputer sendiri}
Jalankan \texttt{bash cek.sh}. Skrip mengompilasi dengan \texttt{-Wall -Werror}, memberi masukan dari \texttt{input/}, lalu mencocokkan keluaran.
\end{BlokGelap}
\end{column}
\end{columns}
\vspace{6pt}
\begin{Catatan}
AI boleh untuk memahami pesan error, tidak untuk meminta solusi utuh. Exit Ticket dikerjakan tanpa bantuan apa pun.
\end{Catatan}
\note{Hands-on tidak dinilai langsung; yang dinilai Exit Ticket dan tugas.}
\end{frame}

\Bagian{5}{Exit Ticket dan tugas}{25 menit terakhir, lalu pekerjaan rumah}
\note{Mulai Exit Ticket tepat waktu.}

\begin{frame}{Exit Ticket 1}
\framesubtitle{Individual, 25 menit, tanpa AI dan internet}
\small
\begin{tabular}{@{}>{\raggedright\arraybackslash}p{60pt}>{\raggedright\arraybackslash}p{560pt}>{\raggedright\arraybackslash}p{80pt}>{\raggedright\arraybackslash}p{150pt}@{}}
\kepala{Soal} & \kepala{Isi} & \kepala{Poin} & \kepala{CPMK}\\
A1 & Tulis program yang menampilkan tiga baris teks yang diberikan & 25 & CPMK0607\\
\barisgenap A2 & Tampilkan teks berisi kutip ganda dan garis miring terbalik & 25 & CPMK0607\\
B1 & Tuliskan keluaran persis dari potongan kode dengan \texttt{endl} dan \texttt{\textbackslash{}n} & 20 & CPMK0608\\
\barisgenap B2 & Temukan dua error pada program dan jelaskan penyebabnya & 20 & CPMK0608\\
B3 & Jelaskan akibat menghapus \texttt{\#include <iostream>} & 10 & CPMK0608\\
\garisdasar
\end{tabular}
\vspace{10pt}

{\small Soal A dinilai dari keluaran yang tepat sama. Soal B dinilai dari ketepatan dan alasan. Pelanggaran aturan membuat nilai Exit Ticket ini 0.}
\note{Soal lengkap dibagikan terpisah dan tidak disimpan di repo publik.}
\end{frame}

\begin{frame}{Tugas 1: Kartu Perkenalan Digital}
\framesubtitle{Dikumpulkan sebelum Pertemuan 2}
\begin{columns}[T]
\begin{column}{0.47\textwidth}
\begin{Blok}{Bagian A, program, 50 poin}
Buat program C++ yang menampilkan kartu perkenalan dengan informasi berikut. Kumpulkan lewat tautan GDB Online (Share) atau kanal yang diumumkan dosen. Ketentuannya di slide berikut.
\end{Blok}
\end{column}
\begin{column}{0.47\textwidth}
\begin{BlokGelap}{Bagian B, penjelasan, 50 poin}
Komentar maksimal 150 kata di atas program: peran \texttt{\#include}, \texttt{main}, dan \texttt{return 0} dengan kalimat kalian sendiri, ditambah satu error yang kalian temui, bunyi pesannya, dan cara memperbaikinya.
\end{BlokGelap}
\end{column}
\end{columns}
\vspace{8pt}
{\small Data di tugas adalah data kalian sendiri. Rubrik lengkap ada di Modul Belajar 1.}
\note{Ingatkan bahwa penjelasan disalin dari AI mudah dikenali karena tidak menyebut pengalaman mereka sendiri.}
\end{frame}

\begin{frame}{Ketentuan Tugas 1}
\framesubtitle{Kartu Perkenalan Digital}
\footnotesize
\begin{tabular}{@{}>{\raggedright\arraybackslash}p{87pt}>{\raggedright\arraybackslash}p{788pt}@{}}
\kepala{No} & \kepala{Ketentuan Bagian A}\\
1 & Header dengan border (karakter \texttt{=} atau \texttt{*})\\
\barisgenap 2 & Nama lengkap\\
3 & NIM\\
\barisgenap 4 & Asal kota\\
5 & Hobi, minimal dua\\
\barisgenap 6 & Cita-cita setelah lulus\\
7 & Footer dengan border; gunakan minimal satu \texttt{\textbackslash{}t} dan satu \texttt{\textbackslash{}"}\\
\barisgenap Bonus & Kreasikan desain kartu, misalnya bingkai ganda atau inisial nama dalam ASCII art.\\
\garisdasar
\end{tabular}
\note{Bacakan ketentuan satu per satu dan tanyakan apakah ada yang belum jelas.}
\end{frame}

\begin{frame}{Rubrik Tugas 1}
\framesubtitle{Empat level, sama di semua instrumen}
\footnotesize
\begin{tabular}{@{}>{\raggedright\arraybackslash}p{166pt}>{\raggedright\arraybackslash}p{44pt}>{\raggedright\arraybackslash}p{166pt}>{\raggedright\arraybackslash}p{156pt}>{\raggedright\arraybackslash}p{146pt}>{\raggedright\arraybackslash}p{146pt}@{}}
\kepala{Kriteria} & \kepala{Poin} & \kepala{Sangat baik 85--100} & \kepala{Baik 70--84} & \kepala{Cukup 55--69} & \kepala{Kurang <55}\\
A1 Program berjalan & 20 & tanpa error dan peringatan & ada peringatan & perlu satu perbaikan & tidak terkompilasi\\
\barisgenap A2 Kelengkapan informasi & 15 & tujuh unsur lengkap & enam unsur & lima unsur & empat atau kurang\\
A3 Format dan kerapian & 15 & rapi, \texttt{\textbackslash{}t} dan \texttt{\textbackslash{}"} tepat & satu escape belum dipakai & ada baris tidak rata & berantakan\\
\barisgenap B1 Penjelasan struktur & 30 & tiga bagian tepat & satu kurang tepat & hanya dua bagian & satu atau disalin\\
B2 Cerita error dan pengujian & 20 & pesan, sebab, perbaikan, uji & pesan tidak dikutip & hanya perbaikan & tidak ada\\
\garisdasar
\end{tabular}
\vspace{8pt}

{\small Nilai kriteria = poin × skor level. Bagian A total 50, Bagian B total 50.}
\note{Rubrik ini sama strukturnya setiap minggu; yang berubah hanya isi kriteria A2, A3, dan B1.}
\end{frame}

\begin{frame}{Sebelum Pertemuan 2}
\framesubtitle{Persiapan}
\begin{columns}[T]
\begin{column}{0.31\textwidth}
\begin{Langkah}{1}{Selesaikan}
Hands-on yang belum lulus \texttt{cek.sh}, terutama latihan tantangan.
\end{Langkah}
\end{column}
\begin{column}{0.31\textwidth}
\begin{Langkah}{2}{Kumpulkan}
Tugas 1 sebelum Pertemuan 2 dimulai.
\end{Langkah}
\end{column}
\begin{column}{0.31\textwidth}
\begin{Langkah}{3}{Baca}
Modul Belajar 1 untuk mengulang, lalu bagian awal modul berikutnya.
\end{Langkah}
\end{column}
\end{columns}
\vspace{10pt}
Pertemuan 2 membahas variabel, tipe data, ekspresi, serta input dengan \texttt{cin}, sejalan dengan Algoritma Pertemuan 2.\note{Tanyakan satu hal yang paling membingungkan hari ini untuk dibuka di review minggu depan.}
\end{frame}

\SlidePenutup{Terima kasih}{Sampai jumpa di Pertemuan 2: variabel, tipe data, dan input}
\note{Slide penutup.}
\end{document}
"
\documentclass[t]{beamer}
\usetheme{ITERA}
% Mode catatan pembicara: aktif bila \catatan didefinisikan sebelum \input.
\ifdefined\catatan
  \setbeameroption{show only notes}
  \setbeamertemplate{note page}{\vspace*{20pt}\hspace*{4pt}\begin{minipage}{900pt}
    {\ITERAKickerFont\fontsize{11pt}{14pt}\selectfont\color{ITERAGoldText}CATATAN PEMBICARA \ \ |\ \ SLIDE \insertframenumber}\par\vspace{4pt}
    {\ITERADisplay\bfseries\fontsize{22pt}{28pt}\selectfont\color{ITERAStructure}\insertframetitle}\par\vspace{12pt}
    \fontsize{15pt}{23pt}\selectfont\insertnote\end{minipage}}
\fi
\tikzset{fase/.style={draw=none,fill=#1,rounded corners=3pt,minimum width=158pt,minimum height=118pt,text width=146pt,align=center}}

\renewcommand\ITERAmeeting{Pertemuan 1}
\begin{document}
\SlideJudul{IF25-11002 PRAKTIKUM PEMROGRAMAN}{Pertemuan 1\\Hello World dan Struktur Program C++}{Program pertama, cara compiler membacanya, dan cara membaca pesan error-nya}{Tim Pengajar}{Tahun 2026. Sejajar dengan Pertemuan 1 Algoritma Pemrograman: Computational Thinking.}
\note{Sapa kelas dan perkenalkan tim pengajar. Hari ini semua orang pulang dengan program yang benar-benar jalan, termasuk yang belum pernah menulis kode.}

\begin{frame}{Dua mata kuliah, satu minggu}
\framesubtitle{Sebelum mulai}
Topik minggu ini sama di dua mata kuliah, dengan peran yang berbeda.
\vspace{10pt}

\begin{columns}[T]
\begin{column}{0.47\textwidth}
\begin{Blok}{Algoritma: Computational Thinking}
Memecah masalah, menulis pseudocode dan flowchart, menelusuri dengan trace table. Minggu ini: computational thinking.
\end{Blok}
\end{column}
\begin{column}{0.47\textwidth}
\begin{BlokGelap}{Praktikum: Hello World}
Menuliskan langkah itu dalam C++, menjalankannya, dan memperbaikinya saat salah. Minggu ini: program pertama.
\end{BlokGelap}
\end{column}
\end{columns}
\note{Tanyakan apa yang sudah dibahas di Algoritma minggu ini. Gunakan jawabannya sebagai jembatan ke praktikum.}
\end{frame}

\begin{frame}{Saat pulang nanti, kalian bisa}
\framesubtitle{Target hari ini}
\begin{columns}[T]
\begin{column}{0.47\textwidth}
\begin{Langkah}{1}{Menerapkan}
Menulis, mengompilasi, dan menjalankan program C++ yang menampilkan keluaran terformat dengan \texttt{cout}, \texttt{endl}, dan escape sequence

{\footnotesize\color{ITERAInkMuted}Sub-CPMK 1.1, CPMK0607}
\end{Langkah}
\end{column}
\begin{column}{0.47\textwidth}
\begin{Langkah}{2}{Menjelaskan}
Menjelaskan peran setiap bagian program C++ dan alur dari kode sumber sampai program berjalan, termasuk membaca pesan error kompilasi

{\footnotesize\color{ITERAInkMuted}Sub-CPMK 1.2, CPMK0608}
\end{Langkah}
\end{column}
\end{columns}
\vspace{8pt}
\begin{Catatan}
Dua kemampuan ini dinilai sama berat di Exit Ticket, tugas, UTS, dan UAS.
\end{Catatan}
\note{Bacakan target dengan pelan. Kata kuncinya menerapkan dan menjelaskan.}
\end{frame}

\begin{frame}{Cara kalian dinilai}
\framesubtitle{Penilaian semester}
\small
\begin{tabular}{@{}>{\raggedright\arraybackslash}p{251pt}>{\raggedright\arraybackslash}p{116pt}>{\raggedright\arraybackslash}p{242pt}>{\raggedright\arraybackslash}p{242pt}@{}}
\kepala{Komponen} & \kepala{Bobot} & \kepala{Bagian A, CPMK0607} & \kepala{Bagian B, CPMK0608}\\
Exit Ticket, 14 kali & 25\% & 12,5 & 12,5\\
\barisgenap Tugas, 12 kali & 25\% & 12,5 & 12,5\\
UTS, pertemuan 8 & 25\% & 12,5 & 12,5\\
\barisgenap UAS, pertemuan 16 & 25\% & 12,5 & 12,5\\
\textbf{Total} & \textbf{100\%} & \textbf{50} & \textbf{50}\\
\garisdasar
\end{tabular}
\vspace{10pt}

Setiap instrumen dibagi rata: Bagian A soal kode, Bagian B soal tracing dan penjelasan. Nilai minimum lulus tiap CPMK adalah 50.\note{Karena setiap instrumen selalu 50:50, kontribusi tiap CPMK pasti 50 poin apa pun bobot komponennya. Sebutkan rubrik empat level yang dipakai di semua tugas.}
\end{frame}

\begin{frame}{Alur 150 menit setiap pertemuan}
\framesubtitle{Siklus yang sama setiap minggu}
\vspace{20pt}
\begin{center}
\begin{tikzpicture}[node distance=14pt]
\node[fase=ITERAGold50,text=ITERAInk] (f0) {{\ITERADisplay\bfseries\fontsize{26pt}{30pt}\selectfont 15'}\\[4pt]{\bfseries Review}\\[2pt]{\footnotesize masalah pembuka}};
\node[fase=ITERAGold50,text=ITERAInk,right=of f0] (f1) {{\ITERADisplay\bfseries\fontsize{26pt}{30pt}\selectfont 25'}\\[4pt]{\bfseries Struggle}\\[2pt]{\footnotesize coba sendiri dulu}};
\draw[-{Stealth[length=7pt]},line width=1.5pt,ITERAStructure] (f0) -- (f1);
\node[fase=ITERAGold,text=ITERAInk,right=of f1] (f2) {{\ITERADisplay\bfseries\fontsize{26pt}{30pt}\selectfont 40'}\\[4pt]{\bfseries Materi}\\[2pt]{\footnotesize tebak lalu buktikan}};
\draw[-{Stealth[length=7pt]},line width=1.5pt,ITERAStructure] (f1) -- (f2);
\node[fase=ITERAGold50,text=ITERAInk,right=of f2] (f3) {{\ITERADisplay\bfseries\fontsize{26pt}{30pt}\selectfont 45'}\\[4pt]{\bfseries Hands-on}\\[2pt]{\footnotesize latihan bertingkat}};
\draw[-{Stealth[length=7pt]},line width=1.5pt,ITERAStructure] (f2) -- (f3);
\node[fase=ITERAGround,text=ITERAOnGround,right=of f3] (f4) {{\ITERADisplay\bfseries\fontsize{26pt}{30pt}\selectfont 25'}\\[4pt]{\bfseries Exit Ticket}\\[2pt]{\footnotesize sendiri, tanpa bantuan}};
\draw[-{Stealth[length=7pt]},line width=1.5pt,ITERAStructure] (f3) -- (f4);
\end{tikzpicture}
\end{center}
\vspace{16pt}
Urutannya sengaja: kalian mencoba dulu sebelum dijelaskan, dan menebak dulu sebelum melihat keluaran.\note{Ini filosofi problem-first dari RPS. Exit Ticket dimulai paling lambat menit ke-125.}
\end{frame}

\Bagian{1}{Masalah pembuka}{Kartu nama di layar}
\note{Masalah pembuka 10 menit.}

\begin{frame}[fragile]{Kartu nama di layar}
\framesubtitle{Masalah minggu ini}
\begin{columns}[T]
\begin{column}{0.55\textwidth}
{\small Bayangkan panitia PKKMB meminta setiap mahasiswa baru menampilkan kartu nama di layar komputer lab. Isinya nama, NIM, dan program studi, dengan garis pembatas di atas dan bawah.

\vspace{6pt}
Kalian tidak boleh mengetik teks itu di Word atau Notepad. Teks harus muncul karena sebuah program menyuruh komputer menampilkannya.\par}
\vspace{6pt}
\begin{Catatan}
\textbf{Diskusi 2 menit.} Kalau komputer hanya mengerti instruksi yang sangat persis, instruksi seperti apa yang perlu kalian tulis? Bagaimana komputer tahu kapan harus pindah baris?
\end{Catatan}
\end{column}
\begin{column}{0.41\textwidth}
\begin{Keluaran}[]
========================
    KARTU MAHASISWA
========================
Nama  : Rani Puspita
NIM   : 126140001
Prodi : Teknik Informatika
========================
\end{Keluaran}
\end{column}
\end{columns}
\note{Tampung beberapa jawaban tanpa menilai benar salah.}
\end{frame}

\begin{frame}{Dari langkah ke instruksi}
\framesubtitle{Sambungan dengan Algoritma}
\begin{columns}[T]
\begin{column}{0.31\textwidth}
\begin{Langkah}{1}{Pecah}
Kartu terdiri dari tujuh baris teks.
\end{Langkah}
\end{column}
\begin{column}{0.31\textwidth}
\begin{Langkah}{2}{Urutkan}
Tampilkan baris pertama, pindah baris, lalu baris kedua, dan seterusnya.
\end{Langkah}
\end{column}
\begin{column}{0.31\textwidth}
\begin{Langkah}{3}{Terjemahkan}
Tulis setiap langkah dalam bahasa yang dimengerti komputer.
\end{Langkah}
\end{column}
\end{columns}
\vspace{10pt}
Langkah 1 dan 2 dilatih di Algoritma. Langkah 3 kita kerjakan hari ini dengan C++.\note{Hubungkan eksplisit ke Algoritma minggu ini.}
\end{frame}

\Bagian{2}{Coba dulu, baru dijelaskan}{25 menit eksperimen di GDB Online}
\note{Struggle 25 menit. Berkeliling, jangan menjelaskan materi dulu.}

\begin{frame}{Misi kalian}
\framesubtitle{Struggle, 25 menit}
\begin{columns}[T]
\begin{column}{0.31\textwidth}
\begin{Langkah}{1}{Buka}
onlinegdb.com, pilih bahasa C++ di kanan atas.
\end{Langkah}
\end{column}
\begin{column}{0.31\textwidth}
\begin{Langkah}{2}{Ubah}
Ganti teks Hello World menjadi Halo, nama saya ..., lalu klik Run.
\end{Langkah}
\end{column}
\begin{column}{0.31\textwidth}
\begin{Langkah}{3}{Kembangkan}
Tampilkan dua baris: nama di baris pertama, NIM di baris kedua.
\end{Langkah}
\end{column}
\end{columns}
\vspace{8pt}
\begin{columns}[T]
\begin{column}{0.47\textwidth}
\begin{Blok}{Boleh}
Bertanya ke teman, mengubah apa saja, sengaja membuat error.
\end{Blok}
\end{column}
\begin{column}{0.47\textwidth}
\begin{BlokAlert}{Belum dulu}
Mencari jawaban jadi di internet atau AI.
\end{BlokAlert}
\end{column}
\end{columns}
\note{Catat kesalahan yang sering muncul untuk dibahas di debrief.}
\end{frame}

\begin{frame}{Apa yang kalian temukan?}
\framesubtitle{Debrief struggle}
\begin{columns}[T]
\begin{column}{0.31\textwidth}
\begin{BlokGelap}{Pertanyaan 1}
Apa yang terjadi kalau satu titik koma dihapus?
\end{BlokGelap}
\end{column}
\begin{column}{0.31\textwidth}
\begin{BlokGelap}{Pertanyaan 2}
Bagaimana cara kalian membuat dua baris?
\end{BlokGelap}
\end{column}
\begin{column}{0.31\textwidth}
\begin{BlokGelap}{Pertanyaan 3}
Baris mana yang menurut kalian tidak boleh dihapus?
\end{BlokGelap}
\end{column}
\end{columns}
\vspace{10pt}
Jawaban kalian akan kita uji di bagian berikutnya.\note{Tulis jawaban di papan; buktikan atau patahkan di bagian materi.}
\end{frame}

\Bagian{3}{Membedah program pertama}{Dari kode sampai keluaran, baris demi baris}
\note{Materi 40 menit. Tunjukkan setiap contoh langsung di GDB Online.}

\begin{frame}{Apa itu program, dan mengapa harus dikompilasi}
\framesubtitle{Kompilasi}
\small
\begin{tabular}{@{}>{\raggedright\arraybackslash}p{194pt}>{\raggedright\arraybackslash}p{681pt}@{}}
\kepala{Tahap} & \kepala{Yang terjadi}\\
1. Kode sumber & Kalian menulis instruksi C++ di file berakhiran \texttt{.cpp}.\\
\barisgenap 2. Kompilasi & Compiler (misalnya \texttt{g++}) memeriksa tata bahasa lalu menerjemahkan ke bahasa mesin.\\
3. Executable & Hasil terjemahan berupa program yang siap dijalankan.\\
\barisgenap 4. Jalankan & Program menampilkan keluaran di layar.\\
\garisdasar
\end{tabular}
\end{frame}

\begin{frame}[fragile]{Membedah Hello World baris demi baris}
\framesubtitle{Hello World}
\begin{columns}[T]
\begin{column}{0.55\textwidth}
\begin{Kode}[]{hello.cpp}
#include <iostream>
using namespace std;

int main() {
    cout << "Hello World!" << endl;
    return 0;
}
\end{Kode}
\end{column}
\begin{column}{0.41\textwidth}
\only<1>{\begin{TebakDulu}
{\ITERAKickerFont\fontsize{10pt}{12pt}\selectfont\color{ITERAAlert}TEBAK DULU}\par\vspace{4pt}
Sebelum dijalankan, tulis keluarannya di kertas.
\end{TebakDulu}
}\only<2>{\KeluaranFile[]{keluaran/P01/k001.txt}
}
\end{column}
\end{columns}
\note<1>{Minta mahasiswa menebak keluaran sebelum membuka slide berikutnya. Tanyakan satu atau dua tebakan yang berbeda, lalu buktikan.}
\end{frame}

\begin{frame}[fragile]{Menyusun keluaran dengan cout, endl, dan \textbackslash{}n}
\framesubtitle{cout}
\begin{columns}[T]
\begin{column}{0.55\textwidth}
\begin{Kode}[]{rangkai.cpp}
cout << "Hai" << "Dunia";
cout << "Hai " << "Dunia" << endl;
cout << "Baris 1" << endl << "Baris 2\n";
cout << "Baris 3\nBaris 4" << endl;
\end{Kode}
\end{column}
\begin{column}{0.41\textwidth}
\only<1>{\begin{TebakDulu}
{\ITERAKickerFont\fontsize{10pt}{12pt}\selectfont\color{ITERAAlert}TEBAK DULU}\par\vspace{4pt}
Sebelum dijalankan, tulis keluarannya di kertas.
\end{TebakDulu}
}\only<2>{\KeluaranFile[]{keluaran/P01/k002.txt}
}
\end{column}
\end{columns}
\note<1>{Minta mahasiswa menebak keluaran sebelum membuka slide berikutnya. Tanyakan satu atau dua tebakan yang berbeda, lalu buktikan.}
\end{frame}

\begin{frame}{Escape sequence untuk karakter khusus}
\framesubtitle{Escape sequence}
\small
\begin{tabular}{@{}>{\raggedright\arraybackslash}p{113pt}>{\raggedright\arraybackslash}p{246pt}>{\raggedright\arraybackslash}p{246pt}>{\raggedright\arraybackslash}p{246pt}@{}}
\kepala{Escape} & \kepala{Arti} & \kepala{Ditulis} & \kepala{Tampil}\\
\texttt{\textbackslash{}n} & baris baru & \texttt{"A\textbackslash{}nB"} & A dan B di dua baris\\
\barisgenap \texttt{\textbackslash{}t} & tab & \texttt{"A\textbackslash{}tB"} & A, jarak tab, B\\
\texttt{\textbackslash{}\textbackslash{}} & garis miring terbalik & \texttt{"C:\textbackslash{}\textbackslash{}data"} & \texttt{C:\textbackslash{}data}\\
\barisgenap \texttt{\textbackslash{}"} & kutip ganda & \texttt{"Dia \textbackslash{}"hai\textbackslash{}""} & \texttt{Dia "hai"}\\
\garisdasar
\end{tabular}
\end{frame}

\begin{frame}[fragile]{Tebak keluarannya}
\framesubtitle{Latihan menjelaskan, 3 menit}
\begin{columns}[T]
\begin{column}{0.56\textwidth}
\begin{Kode}[]{tebak.cpp}
cout << "A" << "B" << endl;
cout << "C\tD\n" << "E";
cout << "\"F\"" << endl;
\end{Kode}
\end{column}
\begin{column}{0.40\textwidth}
\begin{Catatan}
Tulis keluarannya di kertas, karakter demi karakter. Jangan dijalankan dulu.
\end{Catatan}
\end{column}
\end{columns}
\note{Contoh soal Bagian B. Beri waktu 3 menit sebelum slide jawaban.}
\end{frame}

\begin{frame}[fragile]{Tebak keluarannya: jawaban}
\framesubtitle{Latihan menjelaskan}
\begin{columns}[T]
\begin{column}{0.56\textwidth}
\begin{Kode}[]{tebak.cpp}
cout << "A" << "B" << endl;
cout << "C\tD\n" << "E";
cout << "\"F\"" << endl;
\end{Kode}
\end{column}
\begin{column}{0.40\textwidth}
\begin{Keluaran}[]
AB
C       D
E"F"
\end{Keluaran}
\begin{Catatan}
Baris ketiga: \texttt{E} dan \texttt{"F"} menyatu karena tidak ada pindah baris di antaranya. Jarak antara C dan D adalah satu tab.
\end{Catatan}
\end{column}
\end{columns}
\note{Tanyakan jawaban yang paling sering salah, lalu telusuri bersama.}
\end{frame}

\begin{frame}{Error yang sering muncul minggu ini}
\framesubtitle{Pesan diambil langsung dari g++}
\small
\begin{tabular}{@{}>{\raggedright\arraybackslash}p{208pt}>{\raggedright\arraybackslash}p{256pt}>{\raggedright\arraybackslash}p{398pt}@{}}
\kepala{Kesalahan} & \kepala{Contoh} & \kepala{Pesan pertama dari g++}\\
Lupa titik koma & \texttt{cout << "Hai"} & \texttt{error: expected ';' before 'return'}\\
\barisgenap Salah ketik nama & \texttt{count << "Hai";} & \texttt{error: 'count' was not declared in this scope}\\
Lupa \#include & \texttt{tanpa baris 1} & \texttt{error: 'cout' was not declared in this scope}\\
\barisgenap Kutip tidak ditutup & \texttt{cout << "Hai;} & \texttt{error: missing terminating " character}\\
Garis miring tunggal & \texttt{"C:\textbackslash{}data"} & \texttt{warning: unknown escape sequence: '\textbackslash{}d'}\\
\garisdasar
\end{tabular}
\vspace{10pt}

{\small Pesan di tabel ini dihasilkan dengan g++ dan opsi -Wall, sama seperti di cek.sh.}
\note{Minta mahasiswa memotret slide ini.}
\end{frame}

\Bagian{4}{Hands-on bertingkat}{45 menit, dari dasar sampai tantangan}
\note{Mahasiswa membuka folder 02 Hands-on/P1 - Hello World - Hands-on.}

\begin{frame}{Peta hands-on}
\framesubtitle{Folder 02 Hands-on/P1 - Hello World - Hands-on}
\small
\begin{tabular}{@{}>{\raggedright\arraybackslash}p{36pt}>{\raggedright\arraybackslash}p{267pt}>{\raggedright\arraybackslash}p{110pt}>{\raggedright\arraybackslash}p{83pt}>{\raggedright\arraybackslash}p{341pt}@{}}
\kepala{No} & \kepala{File} & \kepala{Tingkat} & \kepala{Waktu} & \kepala{Yang dilatih}\\
1 & \texttt{handson1-biodata.cpp} & Dasar & 10' & \texttt{cout}, \texttt{endl}, tiga baris keluaran\\
\barisgenap 2 & \texttt{handson2-kartu.cpp} & Menengah & 15' & garis pembatas, perataan, \texttt{\textbackslash{}t}, \texttt{\textbackslash{}"}\\
3 & \texttt{handson3-perbaiki.cpp} & Tantangan & 20' & lima kesalahan, perbaiki dan jelaskan\\
\barisgenap + & \texttt{bonus-rumah.cpp} & Pengayaan & sisa & \texttt{\textbackslash{}\textbackslash{}} di dalam ASCII art\\
\garisdasar
\end{tabular}
\vspace{10pt}

Kerjakan berurutan. Cocokkan keluaran dengan folder \texttt{expected/}; latihan yang membaca masukan memakai data di folder \texttt{input/}.\note{Saat berkeliling, minta mahasiswa yang mengerjakan latihan tantangan menjelaskan blok PENJELASAN mereka.}
\end{frame}

\begin{frame}{Cara mengecek dan aturannya}
\framesubtitle{Hands-on}
\begin{columns}[T]
\begin{column}{0.47\textwidth}
\begin{Blok}{Di GDB Online}
Salin isi file latihan, lengkapi TODO, klik Run. Bila program membaca masukan, ketik isi file di \texttt{input/}. Bandingkan dengan \texttt{expected/}.
\end{Blok}
\end{column}
\begin{column}{0.47\textwidth}
\begin{BlokGelap}{Di komputer sendiri}
Jalankan \texttt{bash cek.sh}. Skrip mengompilasi dengan \texttt{-Wall -Werror}, memberi masukan dari \texttt{input/}, lalu mencocokkan keluaran.
\end{BlokGelap}
\end{column}
\end{columns}
\vspace{6pt}
\begin{Catatan}
AI boleh untuk memahami pesan error, tidak untuk meminta solusi utuh. Exit Ticket dikerjakan tanpa bantuan apa pun.
\end{Catatan}
\note{Hands-on tidak dinilai langsung; yang dinilai Exit Ticket dan tugas.}
\end{frame}

\Bagian{5}{Exit Ticket dan tugas}{25 menit terakhir, lalu pekerjaan rumah}
\note{Mulai Exit Ticket tepat waktu.}

\begin{frame}{Exit Ticket 1}
\framesubtitle{Individual, 25 menit, tanpa AI dan internet}
\small
\begin{tabular}{@{}>{\raggedright\arraybackslash}p{60pt}>{\raggedright\arraybackslash}p{560pt}>{\raggedright\arraybackslash}p{80pt}>{\raggedright\arraybackslash}p{150pt}@{}}
\kepala{Soal} & \kepala{Isi} & \kepala{Poin} & \kepala{CPMK}\\
A1 & Tulis program yang menampilkan tiga baris teks yang diberikan & 25 & CPMK0607\\
\barisgenap A2 & Tampilkan teks berisi kutip ganda dan garis miring terbalik & 25 & CPMK0607\\
B1 & Tuliskan keluaran persis dari potongan kode dengan \texttt{endl} dan \texttt{\textbackslash{}n} & 20 & CPMK0608\\
\barisgenap B2 & Temukan dua error pada program dan jelaskan penyebabnya & 20 & CPMK0608\\
B3 & Jelaskan akibat menghapus \texttt{\#include <iostream>} & 10 & CPMK0608\\
\garisdasar
\end{tabular}
\vspace{10pt}

{\small Soal A dinilai dari keluaran yang tepat sama. Soal B dinilai dari ketepatan dan alasan. Pelanggaran aturan membuat nilai Exit Ticket ini 0.}
\note{Soal lengkap dibagikan terpisah dan tidak disimpan di repo publik.}
\end{frame}

\begin{frame}{Tugas 1: Kartu Perkenalan Digital}
\framesubtitle{Dikumpulkan sebelum Pertemuan 2}
\begin{columns}[T]
\begin{column}{0.47\textwidth}
\begin{Blok}{Bagian A, program, 50 poin}
Buat program C++ yang menampilkan kartu perkenalan dengan informasi berikut. Kumpulkan lewat tautan GDB Online (Share) atau kanal yang diumumkan dosen. Ketentuannya di slide berikut.
\end{Blok}
\end{column}
\begin{column}{0.47\textwidth}
\begin{BlokGelap}{Bagian B, penjelasan, 50 poin}
Komentar maksimal 150 kata di atas program: peran \texttt{\#include}, \texttt{main}, dan \texttt{return 0} dengan kalimat kalian sendiri, ditambah satu error yang kalian temui, bunyi pesannya, dan cara memperbaikinya.
\end{BlokGelap}
\end{column}
\end{columns}
\vspace{8pt}
{\small Data di tugas adalah data kalian sendiri. Rubrik lengkap ada di Modul Belajar 1.}
\note{Ingatkan bahwa penjelasan disalin dari AI mudah dikenali karena tidak menyebut pengalaman mereka sendiri.}
\end{frame}

\begin{frame}{Ketentuan Tugas 1}
\framesubtitle{Kartu Perkenalan Digital}
\footnotesize
\begin{tabular}{@{}>{\raggedright\arraybackslash}p{87pt}>{\raggedright\arraybackslash}p{788pt}@{}}
\kepala{No} & \kepala{Ketentuan Bagian A}\\
1 & Header dengan border (karakter \texttt{=} atau \texttt{*})\\
\barisgenap 2 & Nama lengkap\\
3 & NIM\\
\barisgenap 4 & Asal kota\\
5 & Hobi, minimal dua\\
\barisgenap 6 & Cita-cita setelah lulus\\
7 & Footer dengan border; gunakan minimal satu \texttt{\textbackslash{}t} dan satu \texttt{\textbackslash{}"}\\
\barisgenap Bonus & Kreasikan desain kartu, misalnya bingkai ganda atau inisial nama dalam ASCII art.\\
\garisdasar
\end{tabular}
\note{Bacakan ketentuan satu per satu dan tanyakan apakah ada yang belum jelas.}
\end{frame}

\begin{frame}{Rubrik Tugas 1}
\framesubtitle{Empat level, sama di semua instrumen}
\footnotesize
\begin{tabular}{@{}>{\raggedright\arraybackslash}p{166pt}>{\raggedright\arraybackslash}p{44pt}>{\raggedright\arraybackslash}p{166pt}>{\raggedright\arraybackslash}p{156pt}>{\raggedright\arraybackslash}p{146pt}>{\raggedright\arraybackslash}p{146pt}@{}}
\kepala{Kriteria} & \kepala{Poin} & \kepala{Sangat baik 85--100} & \kepala{Baik 70--84} & \kepala{Cukup 55--69} & \kepala{Kurang <55}\\
A1 Program berjalan & 20 & tanpa error dan peringatan & ada peringatan & perlu satu perbaikan & tidak terkompilasi\\
\barisgenap A2 Kelengkapan informasi & 15 & tujuh unsur lengkap & enam unsur & lima unsur & empat atau kurang\\
A3 Format dan kerapian & 15 & rapi, \texttt{\textbackslash{}t} dan \texttt{\textbackslash{}"} tepat & satu escape belum dipakai & ada baris tidak rata & berantakan\\
\barisgenap B1 Penjelasan struktur & 30 & tiga bagian tepat & satu kurang tepat & hanya dua bagian & satu atau disalin\\
B2 Cerita error dan pengujian & 20 & pesan, sebab, perbaikan, uji & pesan tidak dikutip & hanya perbaikan & tidak ada\\
\garisdasar
\end{tabular}
\vspace{8pt}

{\small Nilai kriteria = poin × skor level. Bagian A total 50, Bagian B total 50.}
\note{Rubrik ini sama strukturnya setiap minggu; yang berubah hanya isi kriteria A2, A3, dan B1.}
\end{frame}

\begin{frame}{Sebelum Pertemuan 2}
\framesubtitle{Persiapan}
\begin{columns}[T]
\begin{column}{0.31\textwidth}
\begin{Langkah}{1}{Selesaikan}
Hands-on yang belum lulus \texttt{cek.sh}, terutama latihan tantangan.
\end{Langkah}
\end{column}
\begin{column}{0.31\textwidth}
\begin{Langkah}{2}{Kumpulkan}
Tugas 1 sebelum Pertemuan 2 dimulai.
\end{Langkah}
\end{column}
\begin{column}{0.31\textwidth}
\begin{Langkah}{3}{Baca}
Modul Belajar 1 untuk mengulang, lalu bagian awal modul berikutnya.
\end{Langkah}
\end{column}
\end{columns}
\vspace{10pt}
Pertemuan 2 membahas variabel, tipe data, ekspresi, serta input dengan \texttt{cin}, sejalan dengan Algoritma Pertemuan 2.\note{Tanyakan satu hal yang paling membingungkan hari ini untuk dibuka di review minggu depan.}
\end{frame}

\SlidePenutup{Terima kasih}{Sampai jumpa di Pertemuan 2: variabel, tipe data, dan input}
\note{Slide penutup.}
\end{document}
"
\documentclass[t]{beamer}
\usetheme{ITERA}
% Mode catatan pembicara: aktif bila \catatan didefinisikan sebelum \input.
\ifdefined\catatan
  \setbeameroption{show only notes}
  \setbeamertemplate{note page}{\vspace*{20pt}\hspace*{4pt}\begin{minipage}{900pt}
    {\ITERAKickerFont\fontsize{11pt}{14pt}\selectfont\color{ITERAGoldText}CATATAN PEMBICARA \ \ |\ \ SLIDE \insertframenumber}\par\vspace{4pt}
    {\ITERADisplay\bfseries\fontsize{22pt}{28pt}\selectfont\color{ITERAStructure}\insertframetitle}\par\vspace{12pt}
    \fontsize{15pt}{23pt}\selectfont\insertnote\end{minipage}}
\fi
\tikzset{fase/.style={draw=none,fill=#1,rounded corners=3pt,minimum width=158pt,minimum height=118pt,text width=146pt,align=center}}

\renewcommand\ITERAmeeting{Pertemuan 1}
\begin{document}
\SlideJudul{IF25-11002 PRAKTIKUM PEMROGRAMAN}{Pertemuan 1\\Hello World dan Struktur Program C++}{Program pertama, cara compiler membacanya, dan cara membaca pesan error-nya}{Tim Pengajar}{Tahun 2026. Sejajar dengan Pertemuan 1 Algoritma Pemrograman: Computational Thinking.}
\note{Sapa kelas dan perkenalkan tim pengajar. Hari ini semua orang pulang dengan program yang benar-benar jalan, termasuk yang belum pernah menulis kode.}

\begin{frame}{Dua mata kuliah, satu minggu}
\framesubtitle{Sebelum mulai}
Topik minggu ini sama di dua mata kuliah, dengan peran yang berbeda.
\vspace{10pt}

\begin{columns}[T]
\begin{column}{0.47\textwidth}
\begin{Blok}{Algoritma: Computational Thinking}
Memecah masalah, menulis pseudocode dan flowchart, menelusuri dengan trace table. Minggu ini: computational thinking.
\end{Blok}
\end{column}
\begin{column}{0.47\textwidth}
\begin{BlokGelap}{Praktikum: Hello World}
Menuliskan langkah itu dalam C++, menjalankannya, dan memperbaikinya saat salah. Minggu ini: program pertama.
\end{BlokGelap}
\end{column}
\end{columns}
\note{Tanyakan apa yang sudah dibahas di Algoritma minggu ini. Gunakan jawabannya sebagai jembatan ke praktikum.}
\end{frame}

\begin{frame}{Saat pulang nanti, kalian bisa}
\framesubtitle{Target hari ini}
\begin{columns}[T]
\begin{column}{0.47\textwidth}
\begin{Langkah}{1}{Menerapkan}
Menulis, mengompilasi, dan menjalankan program C++ yang menampilkan keluaran terformat dengan \texttt{cout}, \texttt{endl}, dan escape sequence

{\footnotesize\color{ITERAInkMuted}Sub-CPMK 1.1, CPMK0607}
\end{Langkah}
\end{column}
\begin{column}{0.47\textwidth}
\begin{Langkah}{2}{Menjelaskan}
Menjelaskan peran setiap bagian program C++ dan alur dari kode sumber sampai program berjalan, termasuk membaca pesan error kompilasi

{\footnotesize\color{ITERAInkMuted}Sub-CPMK 1.2, CPMK0608}
\end{Langkah}
\end{column}
\end{columns}
\vspace{8pt}
\begin{Catatan}
Dua kemampuan ini dinilai sama berat di Exit Ticket, tugas, UTS, dan UAS.
\end{Catatan}
\note{Bacakan target dengan pelan. Kata kuncinya menerapkan dan menjelaskan.}
\end{frame}

\begin{frame}{Cara kalian dinilai}
\framesubtitle{Penilaian semester}
\small
\begin{tabular}{@{}>{\raggedright\arraybackslash}p{251pt}>{\raggedright\arraybackslash}p{116pt}>{\raggedright\arraybackslash}p{242pt}>{\raggedright\arraybackslash}p{242pt}@{}}
\kepala{Komponen} & \kepala{Bobot} & \kepala{Bagian A, CPMK0607} & \kepala{Bagian B, CPMK0608}\\
Exit Ticket, 14 kali & 25\% & 12,5 & 12,5\\
\barisgenap Tugas, 12 kali & 25\% & 12,5 & 12,5\\
UTS, pertemuan 8 & 25\% & 12,5 & 12,5\\
\barisgenap UAS, pertemuan 16 & 25\% & 12,5 & 12,5\\
\textbf{Total} & \textbf{100\%} & \textbf{50} & \textbf{50}\\
\garisdasar
\end{tabular}
\vspace{10pt}

Setiap instrumen dibagi rata: Bagian A soal kode, Bagian B soal tracing dan penjelasan. Nilai minimum lulus tiap CPMK adalah 50.\note{Karena setiap instrumen selalu 50:50, kontribusi tiap CPMK pasti 50 poin apa pun bobot komponennya. Sebutkan rubrik empat level yang dipakai di semua tugas.}
\end{frame}

\begin{frame}{Alur 150 menit setiap pertemuan}
\framesubtitle{Siklus yang sama setiap minggu}
\vspace{20pt}
\begin{center}
\begin{tikzpicture}[node distance=14pt]
\node[fase=ITERAGold50,text=ITERAInk] (f0) {{\ITERADisplay\bfseries\fontsize{26pt}{30pt}\selectfont 15'}\\[4pt]{\bfseries Review}\\[2pt]{\footnotesize masalah pembuka}};
\node[fase=ITERAGold50,text=ITERAInk,right=of f0] (f1) {{\ITERADisplay\bfseries\fontsize{26pt}{30pt}\selectfont 25'}\\[4pt]{\bfseries Struggle}\\[2pt]{\footnotesize coba sendiri dulu}};
\draw[-{Stealth[length=7pt]},line width=1.5pt,ITERAStructure] (f0) -- (f1);
\node[fase=ITERAGold,text=ITERAInk,right=of f1] (f2) {{\ITERADisplay\bfseries\fontsize{26pt}{30pt}\selectfont 40'}\\[4pt]{\bfseries Materi}\\[2pt]{\footnotesize tebak lalu buktikan}};
\draw[-{Stealth[length=7pt]},line width=1.5pt,ITERAStructure] (f1) -- (f2);
\node[fase=ITERAGold50,text=ITERAInk,right=of f2] (f3) {{\ITERADisplay\bfseries\fontsize{26pt}{30pt}\selectfont 45'}\\[4pt]{\bfseries Hands-on}\\[2pt]{\footnotesize latihan bertingkat}};
\draw[-{Stealth[length=7pt]},line width=1.5pt,ITERAStructure] (f2) -- (f3);
\node[fase=ITERAGround,text=ITERAOnGround,right=of f3] (f4) {{\ITERADisplay\bfseries\fontsize{26pt}{30pt}\selectfont 25'}\\[4pt]{\bfseries Exit Ticket}\\[2pt]{\footnotesize sendiri, tanpa bantuan}};
\draw[-{Stealth[length=7pt]},line width=1.5pt,ITERAStructure] (f3) -- (f4);
\end{tikzpicture}
\end{center}
\vspace{16pt}
Urutannya sengaja: kalian mencoba dulu sebelum dijelaskan, dan menebak dulu sebelum melihat keluaran.\note{Ini filosofi problem-first dari RPS. Exit Ticket dimulai paling lambat menit ke-125.}
\end{frame}

\Bagian{1}{Masalah pembuka}{Kartu nama di layar}
\note{Masalah pembuka 10 menit.}

\begin{frame}[fragile]{Kartu nama di layar}
\framesubtitle{Masalah minggu ini}
\begin{columns}[T]
\begin{column}{0.55\textwidth}
{\small Bayangkan panitia PKKMB meminta setiap mahasiswa baru menampilkan kartu nama di layar komputer lab. Isinya nama, NIM, dan program studi, dengan garis pembatas di atas dan bawah.

\vspace{6pt}
Kalian tidak boleh mengetik teks itu di Word atau Notepad. Teks harus muncul karena sebuah program menyuruh komputer menampilkannya.\par}
\vspace{6pt}
\begin{Catatan}
\textbf{Diskusi 2 menit.} Kalau komputer hanya mengerti instruksi yang sangat persis, instruksi seperti apa yang perlu kalian tulis? Bagaimana komputer tahu kapan harus pindah baris?
\end{Catatan}
\end{column}
\begin{column}{0.41\textwidth}
\begin{Keluaran}[]
========================
    KARTU MAHASISWA
========================
Nama  : Rani Puspita
NIM   : 126140001
Prodi : Teknik Informatika
========================
\end{Keluaran}
\end{column}
\end{columns}
\note{Tampung beberapa jawaban tanpa menilai benar salah.}
\end{frame}

\begin{frame}{Dari langkah ke instruksi}
\framesubtitle{Sambungan dengan Algoritma}
\begin{columns}[T]
\begin{column}{0.31\textwidth}
\begin{Langkah}{1}{Pecah}
Kartu terdiri dari tujuh baris teks.
\end{Langkah}
\end{column}
\begin{column}{0.31\textwidth}
\begin{Langkah}{2}{Urutkan}
Tampilkan baris pertama, pindah baris, lalu baris kedua, dan seterusnya.
\end{Langkah}
\end{column}
\begin{column}{0.31\textwidth}
\begin{Langkah}{3}{Terjemahkan}
Tulis setiap langkah dalam bahasa yang dimengerti komputer.
\end{Langkah}
\end{column}
\end{columns}
\vspace{10pt}
Langkah 1 dan 2 dilatih di Algoritma. Langkah 3 kita kerjakan hari ini dengan C++.\note{Hubungkan eksplisit ke Algoritma minggu ini.}
\end{frame}

\Bagian{2}{Coba dulu, baru dijelaskan}{25 menit eksperimen di GDB Online}
\note{Struggle 25 menit. Berkeliling, jangan menjelaskan materi dulu.}

\begin{frame}{Misi kalian}
\framesubtitle{Struggle, 25 menit}
\begin{columns}[T]
\begin{column}{0.31\textwidth}
\begin{Langkah}{1}{Buka}
onlinegdb.com, pilih bahasa C++ di kanan atas.
\end{Langkah}
\end{column}
\begin{column}{0.31\textwidth}
\begin{Langkah}{2}{Ubah}
Ganti teks Hello World menjadi Halo, nama saya ..., lalu klik Run.
\end{Langkah}
\end{column}
\begin{column}{0.31\textwidth}
\begin{Langkah}{3}{Kembangkan}
Tampilkan dua baris: nama di baris pertama, NIM di baris kedua.
\end{Langkah}
\end{column}
\end{columns}
\vspace{8pt}
\begin{columns}[T]
\begin{column}{0.47\textwidth}
\begin{Blok}{Boleh}
Bertanya ke teman, mengubah apa saja, sengaja membuat error.
\end{Blok}
\end{column}
\begin{column}{0.47\textwidth}
\begin{BlokAlert}{Belum dulu}
Mencari jawaban jadi di internet atau AI.
\end{BlokAlert}
\end{column}
\end{columns}
\note{Catat kesalahan yang sering muncul untuk dibahas di debrief.}
\end{frame}

\begin{frame}{Apa yang kalian temukan?}
\framesubtitle{Debrief struggle}
\begin{columns}[T]
\begin{column}{0.31\textwidth}
\begin{BlokGelap}{Pertanyaan 1}
Apa yang terjadi kalau satu titik koma dihapus?
\end{BlokGelap}
\end{column}
\begin{column}{0.31\textwidth}
\begin{BlokGelap}{Pertanyaan 2}
Bagaimana cara kalian membuat dua baris?
\end{BlokGelap}
\end{column}
\begin{column}{0.31\textwidth}
\begin{BlokGelap}{Pertanyaan 3}
Baris mana yang menurut kalian tidak boleh dihapus?
\end{BlokGelap}
\end{column}
\end{columns}
\vspace{10pt}
Jawaban kalian akan kita uji di bagian berikutnya.\note{Tulis jawaban di papan; buktikan atau patahkan di bagian materi.}
\end{frame}

\Bagian{3}{Membedah program pertama}{Dari kode sampai keluaran, baris demi baris}
\note{Materi 40 menit. Tunjukkan setiap contoh langsung di GDB Online.}

\begin{frame}{Apa itu program, dan mengapa harus dikompilasi}
\framesubtitle{Kompilasi}
\small
\begin{tabular}{@{}>{\raggedright\arraybackslash}p{194pt}>{\raggedright\arraybackslash}p{681pt}@{}}
\kepala{Tahap} & \kepala{Yang terjadi}\\
1. Kode sumber & Kalian menulis instruksi C++ di file berakhiran \texttt{.cpp}.\\
\barisgenap 2. Kompilasi & Compiler (misalnya \texttt{g++}) memeriksa tata bahasa lalu menerjemahkan ke bahasa mesin.\\
3. Executable & Hasil terjemahan berupa program yang siap dijalankan.\\
\barisgenap 4. Jalankan & Program menampilkan keluaran di layar.\\
\garisdasar
\end{tabular}
\end{frame}

\begin{frame}[fragile]{Membedah Hello World baris demi baris}
\framesubtitle{Hello World}
\begin{columns}[T]
\begin{column}{0.55\textwidth}
\begin{Kode}[]{hello.cpp}
#include <iostream>
using namespace std;

int main() {
    cout << "Hello World!" << endl;
    return 0;
}
\end{Kode}
\end{column}
\begin{column}{0.41\textwidth}
\only<1>{\begin{TebakDulu}
{\ITERAKickerFont\fontsize{10pt}{12pt}\selectfont\color{ITERAAlert}TEBAK DULU}\par\vspace{4pt}
Sebelum dijalankan, tulis keluarannya di kertas.
\end{TebakDulu}
}\only<2>{\KeluaranFile[]{keluaran/P01/k001.txt}
}
\end{column}
\end{columns}
\note<1>{Minta mahasiswa menebak keluaran sebelum membuka slide berikutnya. Tanyakan satu atau dua tebakan yang berbeda, lalu buktikan.}
\end{frame}

\begin{frame}[fragile]{Menyusun keluaran dengan cout, endl, dan \textbackslash{}n}
\framesubtitle{cout}
\begin{columns}[T]
\begin{column}{0.55\textwidth}
\begin{Kode}[]{rangkai.cpp}
cout << "Hai" << "Dunia";
cout << "Hai " << "Dunia" << endl;
cout << "Baris 1" << endl << "Baris 2\n";
cout << "Baris 3\nBaris 4" << endl;
\end{Kode}
\end{column}
\begin{column}{0.41\textwidth}
\only<1>{\begin{TebakDulu}
{\ITERAKickerFont\fontsize{10pt}{12pt}\selectfont\color{ITERAAlert}TEBAK DULU}\par\vspace{4pt}
Sebelum dijalankan, tulis keluarannya di kertas.
\end{TebakDulu}
}\only<2>{\KeluaranFile[]{keluaran/P01/k002.txt}
}
\end{column}
\end{columns}
\note<1>{Minta mahasiswa menebak keluaran sebelum membuka slide berikutnya. Tanyakan satu atau dua tebakan yang berbeda, lalu buktikan.}
\end{frame}

\begin{frame}{Escape sequence untuk karakter khusus}
\framesubtitle{Escape sequence}
\small
\begin{tabular}{@{}>{\raggedright\arraybackslash}p{113pt}>{\raggedright\arraybackslash}p{246pt}>{\raggedright\arraybackslash}p{246pt}>{\raggedright\arraybackslash}p{246pt}@{}}
\kepala{Escape} & \kepala{Arti} & \kepala{Ditulis} & \kepala{Tampil}\\
\texttt{\textbackslash{}n} & baris baru & \texttt{"A\textbackslash{}nB"} & A dan B di dua baris\\
\barisgenap \texttt{\textbackslash{}t} & tab & \texttt{"A\textbackslash{}tB"} & A, jarak tab, B\\
\texttt{\textbackslash{}\textbackslash{}} & garis miring terbalik & \texttt{"C:\textbackslash{}\textbackslash{}data"} & \texttt{C:\textbackslash{}data}\\
\barisgenap \texttt{\textbackslash{}"} & kutip ganda & \texttt{"Dia \textbackslash{}"hai\textbackslash{}""} & \texttt{Dia "hai"}\\
\garisdasar
\end{tabular}
\end{frame}

\begin{frame}[fragile]{Tebak keluarannya}
\framesubtitle{Latihan menjelaskan, 3 menit}
\begin{columns}[T]
\begin{column}{0.56\textwidth}
\begin{Kode}[]{tebak.cpp}
cout << "A" << "B" << endl;
cout << "C\tD\n" << "E";
cout << "\"F\"" << endl;
\end{Kode}
\end{column}
\begin{column}{0.40\textwidth}
\begin{Catatan}
Tulis keluarannya di kertas, karakter demi karakter. Jangan dijalankan dulu.
\end{Catatan}
\end{column}
\end{columns}
\note{Contoh soal Bagian B. Beri waktu 3 menit sebelum slide jawaban.}
\end{frame}

\begin{frame}[fragile]{Tebak keluarannya: jawaban}
\framesubtitle{Latihan menjelaskan}
\begin{columns}[T]
\begin{column}{0.56\textwidth}
\begin{Kode}[]{tebak.cpp}
cout << "A" << "B" << endl;
cout << "C\tD\n" << "E";
cout << "\"F\"" << endl;
\end{Kode}
\end{column}
\begin{column}{0.40\textwidth}
\begin{Keluaran}[]
AB
C       D
E"F"
\end{Keluaran}
\begin{Catatan}
Baris ketiga: \texttt{E} dan \texttt{"F"} menyatu karena tidak ada pindah baris di antaranya. Jarak antara C dan D adalah satu tab.
\end{Catatan}
\end{column}
\end{columns}
\note{Tanyakan jawaban yang paling sering salah, lalu telusuri bersama.}
\end{frame}

\begin{frame}{Error yang sering muncul minggu ini}
\framesubtitle{Pesan diambil langsung dari g++}
\small
\begin{tabular}{@{}>{\raggedright\arraybackslash}p{208pt}>{\raggedright\arraybackslash}p{256pt}>{\raggedright\arraybackslash}p{398pt}@{}}
\kepala{Kesalahan} & \kepala{Contoh} & \kepala{Pesan pertama dari g++}\\
Lupa titik koma & \texttt{cout << "Hai"} & \texttt{error: expected ';' before 'return'}\\
\barisgenap Salah ketik nama & \texttt{count << "Hai";} & \texttt{error: 'count' was not declared in this scope}\\
Lupa \#include & \texttt{tanpa baris 1} & \texttt{error: 'cout' was not declared in this scope}\\
\barisgenap Kutip tidak ditutup & \texttt{cout << "Hai;} & \texttt{error: missing terminating " character}\\
Garis miring tunggal & \texttt{"C:\textbackslash{}data"} & \texttt{warning: unknown escape sequence: '\textbackslash{}d'}\\
\garisdasar
\end{tabular}
\vspace{10pt}

{\small Pesan di tabel ini dihasilkan dengan g++ dan opsi -Wall, sama seperti di cek.sh.}
\note{Minta mahasiswa memotret slide ini.}
\end{frame}

\Bagian{4}{Hands-on bertingkat}{45 menit, dari dasar sampai tantangan}
\note{Mahasiswa membuka folder 02 Hands-on/P1 - Hello World - Hands-on.}

\begin{frame}{Peta hands-on}
\framesubtitle{Folder 02 Hands-on/P1 - Hello World - Hands-on}
\small
\begin{tabular}{@{}>{\raggedright\arraybackslash}p{36pt}>{\raggedright\arraybackslash}p{267pt}>{\raggedright\arraybackslash}p{110pt}>{\raggedright\arraybackslash}p{83pt}>{\raggedright\arraybackslash}p{341pt}@{}}
\kepala{No} & \kepala{File} & \kepala{Tingkat} & \kepala{Waktu} & \kepala{Yang dilatih}\\
1 & \texttt{handson1-biodata.cpp} & Dasar & 10' & \texttt{cout}, \texttt{endl}, tiga baris keluaran\\
\barisgenap 2 & \texttt{handson2-kartu.cpp} & Menengah & 15' & garis pembatas, perataan, \texttt{\textbackslash{}t}, \texttt{\textbackslash{}"}\\
3 & \texttt{handson3-perbaiki.cpp} & Tantangan & 20' & lima kesalahan, perbaiki dan jelaskan\\
\barisgenap + & \texttt{bonus-rumah.cpp} & Pengayaan & sisa & \texttt{\textbackslash{}\textbackslash{}} di dalam ASCII art\\
\garisdasar
\end{tabular}
\vspace{10pt}

Kerjakan berurutan. Cocokkan keluaran dengan folder \texttt{expected/}; latihan yang membaca masukan memakai data di folder \texttt{input/}.\note{Saat berkeliling, minta mahasiswa yang mengerjakan latihan tantangan menjelaskan blok PENJELASAN mereka.}
\end{frame}

\begin{frame}{Cara mengecek dan aturannya}
\framesubtitle{Hands-on}
\begin{columns}[T]
\begin{column}{0.47\textwidth}
\begin{Blok}{Di GDB Online}
Salin isi file latihan, lengkapi TODO, klik Run. Bila program membaca masukan, ketik isi file di \texttt{input/}. Bandingkan dengan \texttt{expected/}.
\end{Blok}
\end{column}
\begin{column}{0.47\textwidth}
\begin{BlokGelap}{Di komputer sendiri}
Jalankan \texttt{bash cek.sh}. Skrip mengompilasi dengan \texttt{-Wall -Werror}, memberi masukan dari \texttt{input/}, lalu mencocokkan keluaran.
\end{BlokGelap}
\end{column}
\end{columns}
\vspace{6pt}
\begin{Catatan}
AI boleh untuk memahami pesan error, tidak untuk meminta solusi utuh. Exit Ticket dikerjakan tanpa bantuan apa pun.
\end{Catatan}
\note{Hands-on tidak dinilai langsung; yang dinilai Exit Ticket dan tugas.}
\end{frame}

\Bagian{5}{Exit Ticket dan tugas}{25 menit terakhir, lalu pekerjaan rumah}
\note{Mulai Exit Ticket tepat waktu.}

\begin{frame}{Exit Ticket 1}
\framesubtitle{Individual, 25 menit, tanpa AI dan internet}
\small
\begin{tabular}{@{}>{\raggedright\arraybackslash}p{60pt}>{\raggedright\arraybackslash}p{560pt}>{\raggedright\arraybackslash}p{80pt}>{\raggedright\arraybackslash}p{150pt}@{}}
\kepala{Soal} & \kepala{Isi} & \kepala{Poin} & \kepala{CPMK}\\
A1 & Tulis program yang menampilkan tiga baris teks yang diberikan & 25 & CPMK0607\\
\barisgenap A2 & Tampilkan teks berisi kutip ganda dan garis miring terbalik & 25 & CPMK0607\\
B1 & Tuliskan keluaran persis dari potongan kode dengan \texttt{endl} dan \texttt{\textbackslash{}n} & 20 & CPMK0608\\
\barisgenap B2 & Temukan dua error pada program dan jelaskan penyebabnya & 20 & CPMK0608\\
B3 & Jelaskan akibat menghapus \texttt{\#include <iostream>} & 10 & CPMK0608\\
\garisdasar
\end{tabular}
\vspace{10pt}

{\small Soal A dinilai dari keluaran yang tepat sama. Soal B dinilai dari ketepatan dan alasan. Pelanggaran aturan membuat nilai Exit Ticket ini 0.}
\note{Soal lengkap dibagikan terpisah dan tidak disimpan di repo publik.}
\end{frame}

\begin{frame}{Tugas 1: Kartu Perkenalan Digital}
\framesubtitle{Dikumpulkan sebelum Pertemuan 2}
\begin{columns}[T]
\begin{column}{0.47\textwidth}
\begin{Blok}{Bagian A, program, 50 poin}
Buat program C++ yang menampilkan kartu perkenalan dengan informasi berikut. Kumpulkan lewat tautan GDB Online (Share) atau kanal yang diumumkan dosen. Ketentuannya di slide berikut.
\end{Blok}
\end{column}
\begin{column}{0.47\textwidth}
\begin{BlokGelap}{Bagian B, penjelasan, 50 poin}
Komentar maksimal 150 kata di atas program: peran \texttt{\#include}, \texttt{main}, dan \texttt{return 0} dengan kalimat kalian sendiri, ditambah satu error yang kalian temui, bunyi pesannya, dan cara memperbaikinya.
\end{BlokGelap}
\end{column}
\end{columns}
\vspace{8pt}
{\small Data di tugas adalah data kalian sendiri. Rubrik lengkap ada di Modul Belajar 1.}
\note{Ingatkan bahwa penjelasan disalin dari AI mudah dikenali karena tidak menyebut pengalaman mereka sendiri.}
\end{frame}

\begin{frame}{Ketentuan Tugas 1}
\framesubtitle{Kartu Perkenalan Digital}
\footnotesize
\begin{tabular}{@{}>{\raggedright\arraybackslash}p{87pt}>{\raggedright\arraybackslash}p{788pt}@{}}
\kepala{No} & \kepala{Ketentuan Bagian A}\\
1 & Header dengan border (karakter \texttt{=} atau \texttt{*})\\
\barisgenap 2 & Nama lengkap\\
3 & NIM\\
\barisgenap 4 & Asal kota\\
5 & Hobi, minimal dua\\
\barisgenap 6 & Cita-cita setelah lulus\\
7 & Footer dengan border; gunakan minimal satu \texttt{\textbackslash{}t} dan satu \texttt{\textbackslash{}"}\\
\barisgenap Bonus & Kreasikan desain kartu, misalnya bingkai ganda atau inisial nama dalam ASCII art.\\
\garisdasar
\end{tabular}
\note{Bacakan ketentuan satu per satu dan tanyakan apakah ada yang belum jelas.}
\end{frame}

\begin{frame}{Rubrik Tugas 1}
\framesubtitle{Empat level, sama di semua instrumen}
\footnotesize
\begin{tabular}{@{}>{\raggedright\arraybackslash}p{166pt}>{\raggedright\arraybackslash}p{44pt}>{\raggedright\arraybackslash}p{166pt}>{\raggedright\arraybackslash}p{156pt}>{\raggedright\arraybackslash}p{146pt}>{\raggedright\arraybackslash}p{146pt}@{}}
\kepala{Kriteria} & \kepala{Poin} & \kepala{Sangat baik 85--100} & \kepala{Baik 70--84} & \kepala{Cukup 55--69} & \kepala{Kurang <55}\\
A1 Program berjalan & 20 & tanpa error dan peringatan & ada peringatan & perlu satu perbaikan & tidak terkompilasi\\
\barisgenap A2 Kelengkapan informasi & 15 & tujuh unsur lengkap & enam unsur & lima unsur & empat atau kurang\\
A3 Format dan kerapian & 15 & rapi, \texttt{\textbackslash{}t} dan \texttt{\textbackslash{}"} tepat & satu escape belum dipakai & ada baris tidak rata & berantakan\\
\barisgenap B1 Penjelasan struktur & 30 & tiga bagian tepat & satu kurang tepat & hanya dua bagian & satu atau disalin\\
B2 Cerita error dan pengujian & 20 & pesan, sebab, perbaikan, uji & pesan tidak dikutip & hanya perbaikan & tidak ada\\
\garisdasar
\end{tabular}
\vspace{8pt}

{\small Nilai kriteria = poin × skor level. Bagian A total 50, Bagian B total 50.}
\note{Rubrik ini sama strukturnya setiap minggu; yang berubah hanya isi kriteria A2, A3, dan B1.}
\end{frame}

\begin{frame}{Sebelum Pertemuan 2}
\framesubtitle{Persiapan}
\begin{columns}[T]
\begin{column}{0.31\textwidth}
\begin{Langkah}{1}{Selesaikan}
Hands-on yang belum lulus \texttt{cek.sh}, terutama latihan tantangan.
\end{Langkah}
\end{column}
\begin{column}{0.31\textwidth}
\begin{Langkah}{2}{Kumpulkan}
Tugas 1 sebelum Pertemuan 2 dimulai.
\end{Langkah}
\end{column}
\begin{column}{0.31\textwidth}
\begin{Langkah}{3}{Baca}
Modul Belajar 1 untuk mengulang, lalu bagian awal modul berikutnya.
\end{Langkah}
\end{column}
\end{columns}
\vspace{10pt}
Pertemuan 2 membahas variabel, tipe data, ekspresi, serta input dengan \texttt{cin}, sejalan dengan Algoritma Pertemuan 2.\note{Tanyakan satu hal yang paling membingungkan hari ini untuk dibuka di review minggu depan.}
\end{frame}

\SlidePenutup{Terima kasih}{Sampai jumpa di Pertemuan 2: variabel, tipe data, dan input}
\note{Slide penutup.}
\end{document}
"
\documentclass[t]{beamer}
\usetheme{ITERA}
% Mode catatan pembicara: aktif bila \catatan didefinisikan sebelum \input.
\ifdefined\catatan
  \setbeameroption{show only notes}
  \setbeamertemplate{note page}{\vspace*{20pt}\hspace*{4pt}\begin{minipage}{900pt}
    {\ITERAKickerFont\fontsize{11pt}{14pt}\selectfont\color{ITERAGoldText}CATATAN PEMBICARA \ \ |\ \ SLIDE \insertframenumber}\par\vspace{4pt}
    {\ITERADisplay\bfseries\fontsize{22pt}{28pt}\selectfont\color{ITERAStructure}\insertframetitle}\par\vspace{12pt}
    \fontsize{15pt}{23pt}\selectfont\insertnote\end{minipage}}
\fi
\tikzset{fase/.style={draw=none,fill=#1,rounded corners=3pt,minimum width=158pt,minimum height=118pt,text width=146pt,align=center}}

\renewcommand\ITERAmeeting{Pertemuan 1}
\begin{document}
\SlideJudul{IF25-11002 PRAKTIKUM PEMROGRAMAN}{Pertemuan 1\\Hello World dan Struktur Program C++}{Program pertama, cara compiler membacanya, dan cara membaca pesan error-nya}{Tim Pengajar}{Tahun 2026. Sejajar dengan Pertemuan 1 Algoritma Pemrograman: Computational Thinking.}
\note{Sapa kelas dan perkenalkan tim pengajar. Hari ini semua orang pulang dengan program yang benar-benar jalan, termasuk yang belum pernah menulis kode.}

\begin{frame}{Dua mata kuliah, satu minggu}
\framesubtitle{Sebelum mulai}
Topik minggu ini sama di dua mata kuliah, dengan peran yang berbeda.
\vspace{10pt}

\begin{columns}[T]
\begin{column}{0.47\textwidth}
\begin{Blok}{Algoritma: Computational Thinking}
Memecah masalah, menulis pseudocode dan flowchart, menelusuri dengan trace table. Minggu ini: computational thinking.
\end{Blok}
\end{column}
\begin{column}{0.47\textwidth}
\begin{BlokGelap}{Praktikum: Hello World}
Menuliskan langkah itu dalam C++, menjalankannya, dan memperbaikinya saat salah. Minggu ini: program pertama.
\end{BlokGelap}
\end{column}
\end{columns}
\note{Tanyakan apa yang sudah dibahas di Algoritma minggu ini. Gunakan jawabannya sebagai jembatan ke praktikum.}
\end{frame}

\begin{frame}{Saat pulang nanti, kalian bisa}
\framesubtitle{Target hari ini}
\begin{columns}[T]
\begin{column}{0.47\textwidth}
\begin{Langkah}{1}{Menerapkan}
Menulis, mengompilasi, dan menjalankan program C++ yang menampilkan keluaran terformat dengan \texttt{cout}, \texttt{endl}, dan escape sequence

{\footnotesize\color{ITERAInkMuted}Sub-CPMK 1.1, CPMK0607}
\end{Langkah}
\end{column}
\begin{column}{0.47\textwidth}
\begin{Langkah}{2}{Menjelaskan}
Menjelaskan peran setiap bagian program C++ dan alur dari kode sumber sampai program berjalan, termasuk membaca pesan error kompilasi

{\footnotesize\color{ITERAInkMuted}Sub-CPMK 1.2, CPMK0608}
\end{Langkah}
\end{column}
\end{columns}
\vspace{8pt}
\begin{Catatan}
Dua kemampuan ini dinilai sama berat di Exit Ticket, tugas, UTS, dan UAS.
\end{Catatan}
\note{Bacakan target dengan pelan. Kata kuncinya menerapkan dan menjelaskan.}
\end{frame}

\begin{frame}{Cara kalian dinilai}
\framesubtitle{Penilaian semester}
\small
\begin{tabular}{@{}>{\raggedright\arraybackslash}p{251pt}>{\raggedright\arraybackslash}p{116pt}>{\raggedright\arraybackslash}p{242pt}>{\raggedright\arraybackslash}p{242pt}@{}}
\kepala{Komponen} & \kepala{Bobot} & \kepala{Bagian A, CPMK0607} & \kepala{Bagian B, CPMK0608}\\
Exit Ticket, 14 kali & 25\% & 12,5 & 12,5\\
\barisgenap Tugas, 12 kali & 25\% & 12,5 & 12,5\\
UTS, pertemuan 8 & 25\% & 12,5 & 12,5\\
\barisgenap UAS, pertemuan 16 & 25\% & 12,5 & 12,5\\
\textbf{Total} & \textbf{100\%} & \textbf{50} & \textbf{50}\\
\garisdasar
\end{tabular}
\vspace{10pt}

Setiap instrumen dibagi rata: Bagian A soal kode, Bagian B soal tracing dan penjelasan. Nilai minimum lulus tiap CPMK adalah 50.\note{Karena setiap instrumen selalu 50:50, kontribusi tiap CPMK pasti 50 poin apa pun bobot komponennya. Sebutkan rubrik empat level yang dipakai di semua tugas.}
\end{frame}

\begin{frame}{Alur 150 menit setiap pertemuan}
\framesubtitle{Siklus yang sama setiap minggu}
\vspace{20pt}
\begin{center}
\begin{tikzpicture}[node distance=14pt]
\node[fase=ITERAGold50,text=ITERAInk] (f0) {{\ITERADisplay\bfseries\fontsize{26pt}{30pt}\selectfont 15'}\\[4pt]{\bfseries Review}\\[2pt]{\footnotesize masalah pembuka}};
\node[fase=ITERAGold50,text=ITERAInk,right=of f0] (f1) {{\ITERADisplay\bfseries\fontsize{26pt}{30pt}\selectfont 25'}\\[4pt]{\bfseries Struggle}\\[2pt]{\footnotesize coba sendiri dulu}};
\draw[-{Stealth[length=7pt]},line width=1.5pt,ITERAStructure] (f0) -- (f1);
\node[fase=ITERAGold,text=ITERAInk,right=of f1] (f2) {{\ITERADisplay\bfseries\fontsize{26pt}{30pt}\selectfont 40'}\\[4pt]{\bfseries Materi}\\[2pt]{\footnotesize tebak lalu buktikan}};
\draw[-{Stealth[length=7pt]},line width=1.5pt,ITERAStructure] (f1) -- (f2);
\node[fase=ITERAGold50,text=ITERAInk,right=of f2] (f3) {{\ITERADisplay\bfseries\fontsize{26pt}{30pt}\selectfont 45'}\\[4pt]{\bfseries Hands-on}\\[2pt]{\footnotesize latihan bertingkat}};
\draw[-{Stealth[length=7pt]},line width=1.5pt,ITERAStructure] (f2) -- (f3);
\node[fase=ITERAGround,text=ITERAOnGround,right=of f3] (f4) {{\ITERADisplay\bfseries\fontsize{26pt}{30pt}\selectfont 25'}\\[4pt]{\bfseries Exit Ticket}\\[2pt]{\footnotesize sendiri, tanpa bantuan}};
\draw[-{Stealth[length=7pt]},line width=1.5pt,ITERAStructure] (f3) -- (f4);
\end{tikzpicture}
\end{center}
\vspace{16pt}
Urutannya sengaja: kalian mencoba dulu sebelum dijelaskan, dan menebak dulu sebelum melihat keluaran.\note{Ini filosofi problem-first dari RPS. Exit Ticket dimulai paling lambat menit ke-125.}
\end{frame}

\Bagian{1}{Masalah pembuka}{Kartu nama di layar}
\note{Masalah pembuka 10 menit.}

\begin{frame}[fragile]{Kartu nama di layar}
\framesubtitle{Masalah minggu ini}
\begin{columns}[T]
\begin{column}{0.55\textwidth}
{\small Bayangkan panitia PKKMB meminta setiap mahasiswa baru menampilkan kartu nama di layar komputer lab. Isinya nama, NIM, dan program studi, dengan garis pembatas di atas dan bawah.

\vspace{6pt}
Kalian tidak boleh mengetik teks itu di Word atau Notepad. Teks harus muncul karena sebuah program menyuruh komputer menampilkannya.\par}
\vspace{6pt}
\begin{Catatan}
\textbf{Diskusi 2 menit.} Kalau komputer hanya mengerti instruksi yang sangat persis, instruksi seperti apa yang perlu kalian tulis? Bagaimana komputer tahu kapan harus pindah baris?
\end{Catatan}
\end{column}
\begin{column}{0.41\textwidth}
\begin{Keluaran}[]
========================
    KARTU MAHASISWA
========================
Nama  : Rani Puspita
NIM   : 126140001
Prodi : Teknik Informatika
========================
\end{Keluaran}
\end{column}
\end{columns}
\note{Tampung beberapa jawaban tanpa menilai benar salah.}
\end{frame}

\begin{frame}{Dari langkah ke instruksi}
\framesubtitle{Sambungan dengan Algoritma}
\begin{columns}[T]
\begin{column}{0.31\textwidth}
\begin{Langkah}{1}{Pecah}
Kartu terdiri dari tujuh baris teks.
\end{Langkah}
\end{column}
\begin{column}{0.31\textwidth}
\begin{Langkah}{2}{Urutkan}
Tampilkan baris pertama, pindah baris, lalu baris kedua, dan seterusnya.
\end{Langkah}
\end{column}
\begin{column}{0.31\textwidth}
\begin{Langkah}{3}{Terjemahkan}
Tulis setiap langkah dalam bahasa yang dimengerti komputer.
\end{Langkah}
\end{column}
\end{columns}
\vspace{10pt}
Langkah 1 dan 2 dilatih di Algoritma. Langkah 3 kita kerjakan hari ini dengan C++.\note{Hubungkan eksplisit ke Algoritma minggu ini.}
\end{frame}

\Bagian{2}{Coba dulu, baru dijelaskan}{25 menit eksperimen di GDB Online}
\note{Struggle 25 menit. Berkeliling, jangan menjelaskan materi dulu.}

\begin{frame}{Misi kalian}
\framesubtitle{Struggle, 25 menit}
\begin{columns}[T]
\begin{column}{0.31\textwidth}
\begin{Langkah}{1}{Buka}
onlinegdb.com, pilih bahasa C++ di kanan atas.
\end{Langkah}
\end{column}
\begin{column}{0.31\textwidth}
\begin{Langkah}{2}{Ubah}
Ganti teks Hello World menjadi Halo, nama saya ..., lalu klik Run.
\end{Langkah}
\end{column}
\begin{column}{0.31\textwidth}
\begin{Langkah}{3}{Kembangkan}
Tampilkan dua baris: nama di baris pertama, NIM di baris kedua.
\end{Langkah}
\end{column}
\end{columns}
\vspace{8pt}
\begin{columns}[T]
\begin{column}{0.47\textwidth}
\begin{Blok}{Boleh}
Bertanya ke teman, mengubah apa saja, sengaja membuat error.
\end{Blok}
\end{column}
\begin{column}{0.47\textwidth}
\begin{BlokAlert}{Belum dulu}
Mencari jawaban jadi di internet atau AI.
\end{BlokAlert}
\end{column}
\end{columns}
\note{Catat kesalahan yang sering muncul untuk dibahas di debrief.}
\end{frame}

\begin{frame}{Apa yang kalian temukan?}
\framesubtitle{Debrief struggle}
\begin{columns}[T]
\begin{column}{0.31\textwidth}
\begin{BlokGelap}{Pertanyaan 1}
Apa yang terjadi kalau satu titik koma dihapus?
\end{BlokGelap}
\end{column}
\begin{column}{0.31\textwidth}
\begin{BlokGelap}{Pertanyaan 2}
Bagaimana cara kalian membuat dua baris?
\end{BlokGelap}
\end{column}
\begin{column}{0.31\textwidth}
\begin{BlokGelap}{Pertanyaan 3}
Baris mana yang menurut kalian tidak boleh dihapus?
\end{BlokGelap}
\end{column}
\end{columns}
\vspace{10pt}
Jawaban kalian akan kita uji di bagian berikutnya.\note{Tulis jawaban di papan; buktikan atau patahkan di bagian materi.}
\end{frame}

\Bagian{3}{Membedah program pertama}{Dari kode sampai keluaran, baris demi baris}
\note{Materi 40 menit. Tunjukkan setiap contoh langsung di GDB Online.}

\begin{frame}{Apa itu program, dan mengapa harus dikompilasi}
\framesubtitle{Kompilasi}
\small
\begin{tabular}{@{}>{\raggedright\arraybackslash}p{194pt}>{\raggedright\arraybackslash}p{681pt}@{}}
\kepala{Tahap} & \kepala{Yang terjadi}\\
1. Kode sumber & Kalian menulis instruksi C++ di file berakhiran \texttt{.cpp}.\\
\barisgenap 2. Kompilasi & Compiler (misalnya \texttt{g++}) memeriksa tata bahasa lalu menerjemahkan ke bahasa mesin.\\
3. Executable & Hasil terjemahan berupa program yang siap dijalankan.\\
\barisgenap 4. Jalankan & Program menampilkan keluaran di layar.\\
\garisdasar
\end{tabular}
\end{frame}

\begin{frame}[fragile]{Membedah Hello World baris demi baris}
\framesubtitle{Hello World}
\begin{columns}[T]
\begin{column}{0.55\textwidth}
\begin{Kode}[]{hello.cpp}
#include <iostream>
using namespace std;

int main() {
    cout << "Hello World!" << endl;
    return 0;
}
\end{Kode}
\end{column}
\begin{column}{0.41\textwidth}
\only<1>{\begin{TebakDulu}
{\ITERAKickerFont\fontsize{10pt}{12pt}\selectfont\color{ITERAAlert}TEBAK DULU}\par\vspace{4pt}
Sebelum dijalankan, tulis keluarannya di kertas.
\end{TebakDulu}
}\only<2>{\KeluaranFile[]{keluaran/P01/k001.txt}
}
\end{column}
\end{columns}
\note<1>{Minta mahasiswa menebak keluaran sebelum membuka slide berikutnya. Tanyakan satu atau dua tebakan yang berbeda, lalu buktikan.}
\end{frame}

\begin{frame}[fragile]{Menyusun keluaran dengan cout, endl, dan \textbackslash{}n}
\framesubtitle{cout}
\begin{columns}[T]
\begin{column}{0.55\textwidth}
\begin{Kode}[]{rangkai.cpp}
cout << "Hai" << "Dunia";
cout << "Hai " << "Dunia" << endl;
cout << "Baris 1" << endl << "Baris 2\n";
cout << "Baris 3\nBaris 4" << endl;
\end{Kode}
\end{column}
\begin{column}{0.41\textwidth}
\only<1>{\begin{TebakDulu}
{\ITERAKickerFont\fontsize{10pt}{12pt}\selectfont\color{ITERAAlert}TEBAK DULU}\par\vspace{4pt}
Sebelum dijalankan, tulis keluarannya di kertas.
\end{TebakDulu}
}\only<2>{\KeluaranFile[]{keluaran/P01/k002.txt}
}
\end{column}
\end{columns}
\note<1>{Minta mahasiswa menebak keluaran sebelum membuka slide berikutnya. Tanyakan satu atau dua tebakan yang berbeda, lalu buktikan.}
\end{frame}

\begin{frame}{Escape sequence untuk karakter khusus}
\framesubtitle{Escape sequence}
\small
\begin{tabular}{@{}>{\raggedright\arraybackslash}p{113pt}>{\raggedright\arraybackslash}p{246pt}>{\raggedright\arraybackslash}p{246pt}>{\raggedright\arraybackslash}p{246pt}@{}}
\kepala{Escape} & \kepala{Arti} & \kepala{Ditulis} & \kepala{Tampil}\\
\texttt{\textbackslash{}n} & baris baru & \texttt{"A\textbackslash{}nB"} & A dan B di dua baris\\
\barisgenap \texttt{\textbackslash{}t} & tab & \texttt{"A\textbackslash{}tB"} & A, jarak tab, B\\
\texttt{\textbackslash{}\textbackslash{}} & garis miring terbalik & \texttt{"C:\textbackslash{}\textbackslash{}data"} & \texttt{C:\textbackslash{}data}\\
\barisgenap \texttt{\textbackslash{}"} & kutip ganda & \texttt{"Dia \textbackslash{}"hai\textbackslash{}""} & \texttt{Dia "hai"}\\
\garisdasar
\end{tabular}
\end{frame}

\begin{frame}[fragile]{Tebak keluarannya}
\framesubtitle{Latihan menjelaskan, 3 menit}
\begin{columns}[T]
\begin{column}{0.56\textwidth}
\begin{Kode}[]{tebak.cpp}
cout << "A" << "B" << endl;
cout << "C\tD\n" << "E";
cout << "\"F\"" << endl;
\end{Kode}
\end{column}
\begin{column}{0.40\textwidth}
\begin{Catatan}
Tulis keluarannya di kertas, karakter demi karakter. Jangan dijalankan dulu.
\end{Catatan}
\end{column}
\end{columns}
\note{Contoh soal Bagian B. Beri waktu 3 menit sebelum slide jawaban.}
\end{frame}

\begin{frame}[fragile]{Tebak keluarannya: jawaban}
\framesubtitle{Latihan menjelaskan}
\begin{columns}[T]
\begin{column}{0.56\textwidth}
\begin{Kode}[]{tebak.cpp}
cout << "A" << "B" << endl;
cout << "C\tD\n" << "E";
cout << "\"F\"" << endl;
\end{Kode}
\end{column}
\begin{column}{0.40\textwidth}
\begin{Keluaran}[]
AB
C       D
E"F"
\end{Keluaran}
\begin{Catatan}
Baris ketiga: \texttt{E} dan \texttt{"F"} menyatu karena tidak ada pindah baris di antaranya. Jarak antara C dan D adalah satu tab.
\end{Catatan}
\end{column}
\end{columns}
\note{Tanyakan jawaban yang paling sering salah, lalu telusuri bersama.}
\end{frame}

\begin{frame}{Error yang sering muncul minggu ini}
\framesubtitle{Pesan diambil langsung dari g++}
\small
\begin{tabular}{@{}>{\raggedright\arraybackslash}p{208pt}>{\raggedright\arraybackslash}p{256pt}>{\raggedright\arraybackslash}p{398pt}@{}}
\kepala{Kesalahan} & \kepala{Contoh} & \kepala{Pesan pertama dari g++}\\
Lupa titik koma & \texttt{cout << "Hai"} & \texttt{error: expected ';' before 'return'}\\
\barisgenap Salah ketik nama & \texttt{count << "Hai";} & \texttt{error: 'count' was not declared in this scope}\\
Lupa \#include & \texttt{tanpa baris 1} & \texttt{error: 'cout' was not declared in this scope}\\
\barisgenap Kutip tidak ditutup & \texttt{cout << "Hai;} & \texttt{error: missing terminating " character}\\
Garis miring tunggal & \texttt{"C:\textbackslash{}data"} & \texttt{warning: unknown escape sequence: '\textbackslash{}d'}\\
\garisdasar
\end{tabular}
\vspace{10pt}

{\small Pesan di tabel ini dihasilkan dengan g++ dan opsi -Wall, sama seperti di cek.sh.}
\note{Minta mahasiswa memotret slide ini.}
\end{frame}

\Bagian{4}{Hands-on bertingkat}{45 menit, dari dasar sampai tantangan}
\note{Mahasiswa membuka folder 02 Hands-on/P1 - Hello World - Hands-on.}

\begin{frame}{Peta hands-on}
\framesubtitle{Folder 02 Hands-on/P1 - Hello World - Hands-on}
\small
\begin{tabular}{@{}>{\raggedright\arraybackslash}p{36pt}>{\raggedright\arraybackslash}p{267pt}>{\raggedright\arraybackslash}p{110pt}>{\raggedright\arraybackslash}p{83pt}>{\raggedright\arraybackslash}p{341pt}@{}}
\kepala{No} & \kepala{File} & \kepala{Tingkat} & \kepala{Waktu} & \kepala{Yang dilatih}\\
1 & \texttt{handson1-biodata.cpp} & Dasar & 10' & \texttt{cout}, \texttt{endl}, tiga baris keluaran\\
\barisgenap 2 & \texttt{handson2-kartu.cpp} & Menengah & 15' & garis pembatas, perataan, \texttt{\textbackslash{}t}, \texttt{\textbackslash{}"}\\
3 & \texttt{handson3-perbaiki.cpp} & Tantangan & 20' & lima kesalahan, perbaiki dan jelaskan\\
\barisgenap + & \texttt{bonus-rumah.cpp} & Pengayaan & sisa & \texttt{\textbackslash{}\textbackslash{}} di dalam ASCII art\\
\garisdasar
\end{tabular}
\vspace{10pt}

Kerjakan berurutan. Cocokkan keluaran dengan folder \texttt{expected/}; latihan yang membaca masukan memakai data di folder \texttt{input/}.\note{Saat berkeliling, minta mahasiswa yang mengerjakan latihan tantangan menjelaskan blok PENJELASAN mereka.}
\end{frame}

\begin{frame}{Cara mengecek dan aturannya}
\framesubtitle{Hands-on}
\begin{columns}[T]
\begin{column}{0.47\textwidth}
\begin{Blok}{Di GDB Online}
Salin isi file latihan, lengkapi TODO, klik Run. Bila program membaca masukan, ketik isi file di \texttt{input/}. Bandingkan dengan \texttt{expected/}.
\end{Blok}
\end{column}
\begin{column}{0.47\textwidth}
\begin{BlokGelap}{Di komputer sendiri}
Jalankan \texttt{bash cek.sh}. Skrip mengompilasi dengan \texttt{-Wall -Werror}, memberi masukan dari \texttt{input/}, lalu mencocokkan keluaran.
\end{BlokGelap}
\end{column}
\end{columns}
\vspace{6pt}
\begin{Catatan}
AI boleh untuk memahami pesan error, tidak untuk meminta solusi utuh. Exit Ticket dikerjakan tanpa bantuan apa pun.
\end{Catatan}
\note{Hands-on tidak dinilai langsung; yang dinilai Exit Ticket dan tugas.}
\end{frame}

\Bagian{5}{Exit Ticket dan tugas}{25 menit terakhir, lalu pekerjaan rumah}
\note{Mulai Exit Ticket tepat waktu.}

\begin{frame}{Exit Ticket 1}
\framesubtitle{Individual, 25 menit, tanpa AI dan internet}
\small
\begin{tabular}{@{}>{\raggedright\arraybackslash}p{60pt}>{\raggedright\arraybackslash}p{560pt}>{\raggedright\arraybackslash}p{80pt}>{\raggedright\arraybackslash}p{150pt}@{}}
\kepala{Soal} & \kepala{Isi} & \kepala{Poin} & \kepala{CPMK}\\
A1 & Tulis program yang menampilkan tiga baris teks yang diberikan & 25 & CPMK0607\\
\barisgenap A2 & Tampilkan teks berisi kutip ganda dan garis miring terbalik & 25 & CPMK0607\\
B1 & Tuliskan keluaran persis dari potongan kode dengan \texttt{endl} dan \texttt{\textbackslash{}n} & 20 & CPMK0608\\
\barisgenap B2 & Temukan dua error pada program dan jelaskan penyebabnya & 20 & CPMK0608\\
B3 & Jelaskan akibat menghapus \texttt{\#include <iostream>} & 10 & CPMK0608\\
\garisdasar
\end{tabular}
\vspace{10pt}

{\small Soal A dinilai dari keluaran yang tepat sama. Soal B dinilai dari ketepatan dan alasan. Pelanggaran aturan membuat nilai Exit Ticket ini 0.}
\note{Soal lengkap dibagikan terpisah dan tidak disimpan di repo publik.}
\end{frame}

\begin{frame}{Tugas 1: Kartu Perkenalan Digital}
\framesubtitle{Dikumpulkan sebelum Pertemuan 2}
\begin{columns}[T]
\begin{column}{0.47\textwidth}
\begin{Blok}{Bagian A, program, 50 poin}
Buat program C++ yang menampilkan kartu perkenalan dengan informasi berikut. Kumpulkan lewat tautan GDB Online (Share) atau kanal yang diumumkan dosen. Ketentuannya di slide berikut.
\end{Blok}
\end{column}
\begin{column}{0.47\textwidth}
\begin{BlokGelap}{Bagian B, penjelasan, 50 poin}
Komentar maksimal 150 kata di atas program: peran \texttt{\#include}, \texttt{main}, dan \texttt{return 0} dengan kalimat kalian sendiri, ditambah satu error yang kalian temui, bunyi pesannya, dan cara memperbaikinya.
\end{BlokGelap}
\end{column}
\end{columns}
\vspace{8pt}
{\small Data di tugas adalah data kalian sendiri. Rubrik lengkap ada di Modul Belajar 1.}
\note{Ingatkan bahwa penjelasan disalin dari AI mudah dikenali karena tidak menyebut pengalaman mereka sendiri.}
\end{frame}

\begin{frame}{Ketentuan Tugas 1}
\framesubtitle{Kartu Perkenalan Digital}
\footnotesize
\begin{tabular}{@{}>{\raggedright\arraybackslash}p{87pt}>{\raggedright\arraybackslash}p{788pt}@{}}
\kepala{No} & \kepala{Ketentuan Bagian A}\\
1 & Header dengan border (karakter \texttt{=} atau \texttt{*})\\
\barisgenap 2 & Nama lengkap\\
3 & NIM\\
\barisgenap 4 & Asal kota\\
5 & Hobi, minimal dua\\
\barisgenap 6 & Cita-cita setelah lulus\\
7 & Footer dengan border; gunakan minimal satu \texttt{\textbackslash{}t} dan satu \texttt{\textbackslash{}"}\\
\barisgenap Bonus & Kreasikan desain kartu, misalnya bingkai ganda atau inisial nama dalam ASCII art.\\
\garisdasar
\end{tabular}
\note{Bacakan ketentuan satu per satu dan tanyakan apakah ada yang belum jelas.}
\end{frame}

\begin{frame}{Rubrik Tugas 1}
\framesubtitle{Empat level, sama di semua instrumen}
\footnotesize
\begin{tabular}{@{}>{\raggedright\arraybackslash}p{166pt}>{\raggedright\arraybackslash}p{44pt}>{\raggedright\arraybackslash}p{166pt}>{\raggedright\arraybackslash}p{156pt}>{\raggedright\arraybackslash}p{146pt}>{\raggedright\arraybackslash}p{146pt}@{}}
\kepala{Kriteria} & \kepala{Poin} & \kepala{Sangat baik 85--100} & \kepala{Baik 70--84} & \kepala{Cukup 55--69} & \kepala{Kurang <55}\\
A1 Program berjalan & 20 & tanpa error dan peringatan & ada peringatan & perlu satu perbaikan & tidak terkompilasi\\
\barisgenap A2 Kelengkapan informasi & 15 & tujuh unsur lengkap & enam unsur & lima unsur & empat atau kurang\\
A3 Format dan kerapian & 15 & rapi, \texttt{\textbackslash{}t} dan \texttt{\textbackslash{}"} tepat & satu escape belum dipakai & ada baris tidak rata & berantakan\\
\barisgenap B1 Penjelasan struktur & 30 & tiga bagian tepat & satu kurang tepat & hanya dua bagian & satu atau disalin\\
B2 Cerita error dan pengujian & 20 & pesan, sebab, perbaikan, uji & pesan tidak dikutip & hanya perbaikan & tidak ada\\
\garisdasar
\end{tabular}
\vspace{8pt}

{\small Nilai kriteria = poin × skor level. Bagian A total 50, Bagian B total 50.}
\note{Rubrik ini sama strukturnya setiap minggu; yang berubah hanya isi kriteria A2, A3, dan B1.}
\end{frame}

\begin{frame}{Sebelum Pertemuan 2}
\framesubtitle{Persiapan}
\begin{columns}[T]
\begin{column}{0.31\textwidth}
\begin{Langkah}{1}{Selesaikan}
Hands-on yang belum lulus \texttt{cek.sh}, terutama latihan tantangan.
\end{Langkah}
\end{column}
\begin{column}{0.31\textwidth}
\begin{Langkah}{2}{Kumpulkan}
Tugas 1 sebelum Pertemuan 2 dimulai.
\end{Langkah}
\end{column}
\begin{column}{0.31\textwidth}
\begin{Langkah}{3}{Baca}
Modul Belajar 1 untuk mengulang, lalu bagian awal modul berikutnya.
\end{Langkah}
\end{column}
\end{columns}
\vspace{10pt}
Pertemuan 2 membahas variabel, tipe data, ekspresi, serta input dengan \texttt{cin}, sejalan dengan Algoritma Pertemuan 2.\note{Tanyakan satu hal yang paling membingungkan hari ini untuk dibuka di review minggu depan.}
\end{frame}

\SlidePenutup{Terima kasih}{Sampai jumpa di Pertemuan 2: variabel, tipe data, dan input}
\note{Slide penutup.}
\end{document}
